\documentclass[12pt,a4paper]{article}

% ── Encoding & Language ──────────────────────────────────────────────
\usepackage[UTF8]{ctex}
\usepackage[utf8]{inputenc}
\usepackage[T1]{fontenc}

% ── Math & AMS ───────────────────────────────────────────────────────
\usepackage{amsmath,amssymb,amsthm}
\usepackage{mathtools}
\usepackage{mathrsfs}

% ── Geometry & Diagrams ──────────────────────────────────────────────
\usepackage{tikz-cd}
\usepackage{graphicx}

% ── Formatting ───────────────────────────────────────────────────────
\usepackage{geometry}
\geometry{margin=2.5cm}
\usepackage{hyperref}
\usepackage{enumitem}
\usepackage{booktabs}
\usepackage{bm}

% ── Theorem environments ─────────────────────────────────────────────
\theoremstyle{definition}
\newtheorem{definition}{定义}[section]
\newtheorem{theorem}{定理}[section]
\newtheorem{proposition}{命题}[section]
\newtheorem{corollary}{推论}[section]
\newtheorem{lemma}{引理}[section]
\newtheorem{remark}{注记}[section]
\newtheorem{example}{例}[section]

\theoremstyle{plain}
\newtheorem*{definition*}{Definition}
\newtheorem*{theorem*}{Theorem}
\newtheorem*{proposition*}{Proposition}
\newtheorem*{corollary*}{Corollary}
\newtheorem*{lemma*}{Lemma}
\newtheorem*{remark*}{Remark}

% ── Custom commands ──────────────────────────────────────────────────
\newcommand{\R}{\mathbb{R}}
\newcommand{\C}{\mathbb{C}}
\newcommand{\Z}{\mathbb{Z}}
\newcommand{\N}{\mathbb{N}}
\newcommand{\g}{\mathfrak{g}}
\newcommand{\lie}[1]{\mathfrak{#1}}
\newcommand{\dif}{\mathrm{d}}
\newcommand{\curv}{\Omega}
\newcommand{\conn}{\omega}
\newcommand{\bun}{\mathcal{P}}
\newcommand{\base}{\mathcal{M}}
\newcommand{\Ad}{\mathrm{Ad}}
\newcommand{\ad}{\mathrm{ad}}
\newcommand{\Hor}{\mathrm{Hor}}
\newcommand{\Ver}{\mathrm{Ver}}
\newcommand{\Hol}{\mathrm{Hol}}
\newcommand{\Tr}{\mathrm{Tr}}
\newcommand{\vol}{\mathrm{vol}}
\newcommand{\cercis}{\mathcal{C}}
\newcommand{\Iso}{\mathrm{Iso}}

\title{
    {\Huge \textbf{SCX 规范理论的纤维丛几何形式化}}\\[0.5em]
    {\Large Formalization of SCX Gauge Theory as Fiber Bundle Geometry}\\[1em]
    {\normalsize 规范自由度、曲率障碍与 MILP 库仑规范}\\[0.3em]
    {\normalsize Gauge Freedom, Curvature Obstruction, and MILP Coulomb Gauge}
}

\author{}
\date{\today}

\begin{document}
\maketitle

% ══════════════════════════════════════════════════════════════════════
\begin{abstract}
\noindent
{\bf 中文摘要.}
本文从微分几何的主纤维丛理论出发，严格形式化了 SCX (Self-Consistent eXpert) 框架中的规范理论结构。
我们将每个专家的平移群视为一个 Abel 规范群 $G_m \cong (\R^d, +)$，
全体专家的直积 $G = \prod_{m=1}^M G_m$ 构成一个主 $G$-丛 $\bun \to \base$，
其中底流形 $\base$ 是所有物理上可区分的专家构型空间（规范轨道）。
通过定义联络 $1$-形式 $\conn$ 及其曲率 $2$-形式 $\curv = \dif \conn$，
我们证明了：{\bf (i)} $\curv = 0$ 当且仅当存在全局规范固定（所有专家共享同一规范）；
{\bf (ii)} $\curv \neq 0$ 时，不存在同时对齐所有成对专家比较的单一规范选择 —— 这便是``势能面不齐''的几何根源；
{\bf (iii)} Cercis 分数是规范不变的几何可观测量，类似于 Yang-Mills 理论中的 $\Tr(F_{\mu\nu}F^{\mu\nu})$；
{\bf (iv)} MILP 规范固定是离散情形的 Coulomb 规范，满足 $\sum_m g_m = 0$ 的约束最小化联络的 $L^2$ 范数。
该框架为多专家系统的全局可比较性提供了严格的拓扑判别准则。

\vspace{0.8em}
\noindent
{\bf English Abstract.}
This paper rigorously formalizes the gauge-theoretic structure of the SCX (Self-Consistent eXpert) framework
using the language of principal fiber bundles in differential geometry.
We identify the translation group of each expert as an abelian gauge group $G_m \cong (\R^d, +)$,
and the direct product $G = \prod_{m=1}^M G_m$ as the structure group of a principal $G$-bundle
$\bun \to \base$, where the base manifold $\base$ is the space of all physically distinguishable
expert configurations (gauge orbits).
By introducing a connection $1$-form $\conn$ and its curvature $2$-form $\curv = \dif \conn$, we prove:
{\bf (i)} $\curv = 0$ iff global gauge-fixing is possible (all experts share the same gauge);
{\bf (ii)} when $\curv \neq 0$, no single gauge choice can simultaneously align all pairwise expert
comparisons — this is the geometric origin of ``potential energy surface misalignment'' (势能面不齐);
{\bf (iii)} the Cercis Score is a gauge-invariant geometric observable, analogous to
$\Tr(F_{\mu\nu}F^{\mu\nu})$ in Yang-Mills theory;
{\bf (iv)} MILP gauge-fixing is a discrete analog of Coulomb gauge selection, satisfying the
constraint $\sum_m g_m = 0$ and minimizing the $L^2$ norm of the connection.
This framework provides a rigorous topological criterion for global comparability in multi-expert systems.
\end{abstract}

\tableofcontents
\newpage

% ══════════════════════════════════════════════════════════════════════
\section{引言 / Introduction}
% ══════════════════════════════════════════════════════════════════════

\subsection{问题的提出 / Motivation}

在多专家系统中，每个专家为其观测空间赋予一个内部坐标系（规范），
不同专家之间坐标系的变换构成了一个规范群作用。
当多个专家对同一物理系统进行预测或评估时，如何判断不同专家的输出
是否属于同一规范类、是否可以进行有意义的比较，是一个核心问题。

SCX (Self-Consistent eXpert) 框架将每个专家建模为在一个 $d$ 维构型空间上的
``地图绘制者''（cartographer）。专家 $m$ 的输出 $x_m \in \R^d$ 是该构型空间中
一个点的坐标表示，依赖于专家选择的规范 $g_m \in G_m$。

两个核心问题随之产生：
\begin{enumerate}[label=(\roman*)]
    \item {\bf 全局规范固定问题 (Global Gauge-Fixing Problem):}
        是否存在一个全局的规范选择，使得所有专家的坐标表示可以直接比较？
    \item {\bf 规范不变量的构造 (Gauge-Invariant Observables):}
        能否构造不依赖规范选择的物理量，以客观地度量多专家系统的一致性？
\end{enumerate}

本文的核心贡献在于：将上述问题映射到微分几何中主纤维丛上的联络与曲率理论，
从而利用曲率作为障碍类的拓扑性质给出严格的回答。

\medskip
\noindent
In multi-expert systems, each expert assigns an internal coordinate system (gauge)
to its observation space. The transformations between different experts' coordinate
systems constitute a gauge group action. When multiple experts predict or evaluate
the same physical system, a central question is how to determine whether outputs
from different experts belong to the same gauge class and can be meaningfully compared.

The SCX (Self-Consistent eXpert) framework models each expert as a ``cartographer''
on a $d$-dimensional configuration space. The output $x_m \in \R^d$ of expert $m$
is a coordinate representation of a point in this configuration space, depending on
the gauge $g_m \in G_m$ chosen by that expert.

Two fundamental questions arise:
\begin{enumerate}[label=(\roman*)]
    \item {\bf Global Gauge-Fixing Problem:}
        Does there exist a global gauge choice such that all experts' coordinate
        representations are directly comparable?
    \item {\bf Construction of Gauge-Invariant Observables:}
        Can we construct physical quantities independent of gauge choices to
        objectively measure the consistency of a multi-expert system?
\end{enumerate}

The central contribution of this paper is to map these questions onto the theory
of connections and curvature on principal fiber bundles in differential geometry,
thereby providing rigorous answers using the topological nature of curvature as
an obstruction class.

\subsection{相关工作与动机 / Related Work and Motivation}

规范理论在物理学中有深厚的根基，从 Weyl 的标度规范理论~\cite{weyl1918}
到 Yang-Mills 的非 Abel 规范理论~\cite{yangmills1954}，再到现代粒子物理标准模型。
在数学上，Atiyah 与 Singer、Donaldson、Witten 等人的工作深刻揭示了规范理论
与拓扑学之间的联系。

然而，将规范理论形式化地应用于机器学习中的多专家系统比对问题，
尚属首次。本文借鉴了以下数学结构：
\begin{itemize}
    \item 主纤维丛 (principal fiber bundles)~\cite{kobayashi1963}：描述规范自由度；
    \item Ehresmann 联络 (Ehresmann connection)~\cite{ehresmann1950}：描述不同纤维之间的``平行移动''；
    \item 曲率作为障碍 (curvature as obstruction)：度量平行移动的路径依赖性；
    \item 和乐群 (holonomy group)：分类非平凡规范结构。
\end{itemize}

\medskip
\noindent
Gauge theory has deep roots in physics, from Weyl's scale gauge theory~\cite{weyl1918}
to Yang-Mills non-abelian gauge theory~\cite{yangmills1954}, and to the modern Standard Model
of particle physics. Mathematically, the works of Atiyah--Singer, Donaldson, and Witten have
profoundly revealed the connections between gauge theory and topology.

However, the formal application of gauge theory to the problem of comparing multi-expert
systems in machine learning is, to our knowledge, novel. This paper draws on the following
mathematical structures:
\begin{itemize}
    \item Principal fiber bundles~\cite{kobayashi1963}: describing gauge degrees of freedom;
    \item Ehresmann connections~\cite{ehresmann1950}: describing ``parallel transport'' between fibers;
    \item Curvature as obstruction: measuring the path-dependence of parallel transport;
    \item Holonomy groups: classifying non-trivial gauge structures.
\end{itemize}

% ══════════════════════════════════════════════════════════════════════
\section{数学预备 / Mathematical Preliminaries}
% ══════════════════════════════════════════════════════════════════════

\subsection{主纤维丛 / Principal Fiber Bundles}

\begin{definition}[主纤维丛]
设 $G$ 为 Lie 群，$\bun$ 为光滑流形，$\base$ 为光滑流形，
$\pi: \bun \to \base$ 为光滑满浸没映射。
若 $G$ 在 $\bun$ 上有一个自由右作用
\[
    R: \bun \times G \to \bun, \quad (p, g) \mapsto R_g(p) = p \cdot g
\]
且满足：
\begin{enumerate}[label=(\arabic*)]
    \item $\pi(p \cdot g) = \pi(p)$ 对所有 $p \in \bun, g \in G$ 成立；
    \item（局部平凡性）对每个 $x \in \base$，存在开邻域 $U \subset \base$ 及微分同胚
        $\phi: \pi^{-1}(U) \to U \times G$，使得
        $\phi(p) = (\pi(p), \psi(p))$，
        其中 $\psi: \pi^{-1}(U) \to G$ 满足 $\psi(p \cdot g) = \psi(p) g$。
\end{enumerate}
则称 $(\bun, \base, \pi, G)$ 为以 $G$ 为结构群、$\base$ 为底流形的主纤维丛
(principal fiber bundle，简称主丛)。
\end{definition}

\noindent
Let $G$ be a Lie group, $\bun$ a smooth manifold, $\base$ a smooth manifold,
and $\pi: \bun \to \base$ a smooth surjective submersion.
If $G$ acts freely from the right on $\bun$ via
$R: \bun \times G \to \bun$, $(p, g) \mapsto R_g(p) = p \cdot g$,
satisfying:
\begin{enumerate}[label=(\arabic*)]
    \item $\pi(p \cdot g) = \pi(p)$ for all $p \in \bun, g \in G$;
    \item (Local triviality) for each $x \in \base$, there exists an open neighborhood $U \subset \base$
        and a diffeomorphism $\phi: \pi^{-1}(U) \to U \times G$ such that
        $\phi(p) = (\pi(p), \psi(p))$,
        where $\psi: \pi^{-1}(U) \to G$ satisfies $\psi(p \cdot g) = \psi(p) g$,
\end{enumerate}
then $(\bun, \base, \pi, G)$ is called a principal fiber bundle with structure group $G$
and base manifold $\base$.

\begin{definition}[局部截面、规范选择]
光滑映射 $s: U \to \bun$（其中 $U \subset \base$ 为开集）称为局部截面 (local section)，
若 $\pi \circ s = \id_U$。
一个局部截面在物理上对应于一个局部的规范选择 (local gauge choice)。
\end{definition}

\noindent
A smooth map $s: U \to \bun$ (where $U \subset \base$ is open) is called a local section
if $\pi \circ s = \id_U$. A local section corresponds physically to a local gauge choice.

\subsection{切空间与竖直子丛 / Tangent Space and Vertical Subbundle}

\begin{definition}[竖直子空间]
在 $p \in \bun$ 处，定义竖直子空间
\[
    V_p \bun := \ker(\dif \pi_p) \subset T_p \bun.
\]
$V_p \bun$ 中向量沿纤维方向，称为竖直向量 (vertical vectors)。
所有 $V_p\bun$ 构成 $T\bun$ 的可积子丛 $V\bun$。
\end{definition}

\noindent
At $p \in \bun$, define the vertical subspace
$V_p \bun := \ker(\dif \pi_p) \subset T_p \bun$.
Vectors in $V_p \bun$ are along the fiber direction and are called vertical vectors.
All $V_p\bun$ form an integrable subbundle $V\bun$ of $T\bun$.

\begin{definition}[基本向量场]
对每个 $\xi \in \g = \lie{G}$（$G$ 的 Lie 代数），定义 $\bun$ 上的基本向量场
(fundamental vector field) $\xi^\sharp$：
\[
    \xi^\sharp_p := \left.\frac{\dif}{\dif t}\right|_{t=0} (p \cdot \exp(t\xi)).
\]
映射 $\xi \mapsto \xi^\sharp_p$ 给出 Lie 代数同构 $\g \xrightarrow{\cong} V_p\bun$。
\end{definition}

\noindent
For each $\xi \in \g = \lie{G}$ (the Lie algebra of $G$), define the fundamental vector field
$\xi^\sharp$ on $\bun$:
$\xi^\sharp_p := \left.\frac{\dif}{\dif t}\right|_{t=0} (p \cdot \exp(t\xi))$.
The map $\xi \mapsto \xi^\sharp_p$ gives a Lie algebra isomorphism $\g \xrightarrow{\cong} V_p\bun$.

\subsection{Ehresmann 联络 / Ehresmann Connection}

\begin{definition}[Ehresmann 联络]
$\bun$ 上的 Ehresmann 联络是 $T\bun$ 的一个光滑子丛 $H\bun$，满足：
\begin{enumerate}[label=(\arabic*)]
    \item $T_p \bun = V_p \bun \oplus H_p \bun$ 对所有 $p \in \bun$ 成立；
    \item（$G$-等变性）$H_{p \cdot g} \bun = (R_g)_* H_p \bun$ 对所有 $p \in \bun, g \in G$ 成立。
\end{enumerate}
$H_p\bun$ 中的向量称为水平向量 (horizontal vectors)。
\end{definition}

\noindent
An Ehresmann connection on $\bun$ is a smooth subbundle $H\bun$ of $T\bun$ satisfying:
\begin{enumerate}[label=(\arabic*)]
    \item $T_p \bun = V_p \bun \oplus H_p \bun$ for all $p \in \bun$;
    \item ($G$-equivariance) $H_{p \cdot g} \bun = (R_g)_* H_p \bun$ for all $p \in \bun, g \in G$.
\end{enumerate}
Vectors in $H_p\bun$ are called horizontal vectors.

\begin{definition}[联络1-形式 / Connection 1-Form]
Ehresmann 联络等价地由 $G$-等变的 $\g$-值 $1$-形式 $\conn \in \Omega^1(\bun, \g)$ 给出，满足：
\[
    \conn(\xi^\sharp) = \xi, \quad \forall \xi \in \g,
\]
及 $G$-等变条件：
\[
    R_g^* \conn = \Ad_{g^{-1}} \circ \conn.
\]
此时 $H_p\bun = \ker(\conn_p)$。
\end{definition}

\noindent
An Ehresmann connection is equivalently given by a $G$-equivariant $\g$-valued $1$-form
$\conn \in \Omega^1(\bun, \g)$, satisfying:
$\conn(\xi^\sharp) = \xi$ for all $\xi \in \g$,
and the $G$-equivariance condition:
$R_g^* \conn = \Ad_{g^{-1}} \circ \conn$.
Then $H_p\bun = \ker(\conn_p)$.

\subsection{曲率 / Curvature}

\begin{definition}[曲率2-形式]
联络 $\conn$ 的曲率 $2$-形式 $\curv \in \Omega^2(\bun, \g)$ 定义为 $\conn$ 的外共变导数：
\[
    \curv := \dif^{\conn} \conn := \dif \conn + \frac{1}{2}[\conn, \conn],
\]
其中 $[\conn, \conn](X, Y) := [\conn(X), \conn(Y)]$。
当 $G$ 为 Abel 群时，$[\conn, \conn] = 0$，故 $\curv = \dif \conn$。
\end{definition}

\noindent
The curvature $2$-form $\curv \in \Omega^2(\bun, \g)$ of connection $\conn$ is defined as
the exterior covariant derivative of $\conn$:
$\curv := \dif^{\conn} \conn := \dif \conn + \frac{1}{2}[\conn, \conn]$,
where $[\conn, \conn](X, Y) := [\conn(X), \conn(Y)]$.
When $G$ is abelian, $[\conn, \conn] = 0$, so $\curv = \dif \conn$.

曲率满足 Bianchi 恒等式：
\[
    \dif^{\conn} \curv = \dif \curv + [\conn, \curv] = 0.
\]

The curvature satisfies the Bianchi identity:
$\dif^{\conn} \curv = \dif \curv + [\conn, \curv] = 0$.

\subsection{拉回与规范势 / Pullback and Gauge Potential}

\begin{definition}[规范势]
给定局部截面 $s: U \to \bun$（即局部规范选择），定义局部规范势
(local gauge potential) $A^s \in \Omega^1(U, \g)$ 及局部场强
$F^s \in \Omega^2(U, \g)$：
\[
    A^s := s^* \conn, \qquad F^s := s^* \curv.
\]
在 $U$ 上 $F^s = \dif A^s + \frac{1}{2}[A^s, A^s]$（Abel 情形为 $F^s = \dif A^s$）。
\end{definition}

\noindent
Given a local section $s: U \to \bun$ (i.e., a local gauge choice), define the local gauge potential
$A^s \in \Omega^1(U, \g)$ and local field strength $F^s \in \Omega^2(U, \g)$:
$A^s := s^* \conn$, $F^s := s^* \curv$.
On $U$, $F^s = \dif A^s + \frac{1}{2}[A^s, A^s]$ (for abelian case, $F^s = \dif A^s$).

\begin{definition}[规范变换]
若 $s, s': U \to \bun$ 为两个局部截面，则存在光滑映射 $g: U \to G$（规范变换函数）
使得 $s'(x) = s(x) \cdot g(x)$。
此时规范势的变换规则为：
\[
    A^{s'} = \Ad_{g^{-1}} \circ A^s + g^* \theta,
\]
其中 $\theta \in \Omega^1(G, \g)$ 是 $G$ 上的 Maurer-Cartan 形式。
对于 Abel 群，$G_m \cong \R^d$ 且 $\theta = \dif g$（视为 $\R^d$ 上标准平移），故：
\[
    A^{s'} = A^s - \dif g.
\]
\end{definition}

\noindent
If $s, s': U \to \bun$ are two local sections, there exists a smooth map $g: U \to G$
(the gauge transformation function) such that $s'(x) = s(x) \cdot g(x)$.
The gauge potential transforms as:
$A^{s'} = \Ad_{g^{-1}} \circ A^s + g^* \theta$,
where $\theta \in \Omega^1(G, \g)$ is the Maurer-Cartan form on $G$.
For abelian groups, $G_m \cong \R^d$ and $\theta = \dif g$, so:
$A^{s'} = A^s - \dif g$.

\subsection{和乐群 / Holonomy Group}

\begin{definition}[和乐]
设 $\gamma: [0,1] \to \base$ 为以 $x_0 \in \base$ 为基点的分段光滑闭曲线。
沿 $\gamma$ 的平行移动定义了纤维 $\pi^{-1}(x_0)$ 上的自同构。
所有这样的自同构构成的群称为（约化）和乐群 $\Hol_{x_0}(\conn) \subset G$。
\end{definition}

\noindent
Let $\gamma: [0,1] \to \base$ be a piecewise smooth closed curve based at $x_0 \in \base$.
Parallel transport along $\gamma$ defines an automorphism of the fiber $\pi^{-1}(x_0)$.
The group of all such automorphisms is the (restricted) holonomy group
$\Hol_{x_0}(\conn) \subset G$.

\begin{theorem}[Ambrose-Singer, 平坦联络的特征]
$\curv = 0$ 当且仅当 $\Hol_{x_0}(\conn)$ 是离散的（特别地，对单连通底流形，$\Hol_{x_0}(\conn) = \{e\}$）。
此时联络称为{\bf 平坦联络 (flat connection)}，且局部上存在规范
使得 $A^s = 0$。
\end{theorem}

\noindent
$\curv = 0$ if and only if $\Hol_{x_0}(\conn)$ is discrete (in particular,
for simply connected base manifolds, $\Hol_{x_0}(\conn) = \{e\}$).
Such a connection is called {\bf flat}, and locally there exists a gauge such that $A^s = 0$.

% ══════════════════════════════════════════════════════════════════════
\section{SCX 规范群 / The SCX Gauge Group}
% ══════════════════════════════════════════════════════════════════════

\subsection{单个专家的平移群 / Translation Group of a Single Expert}

在 SCX 框架中，每个专家 $m \in \{1, 2, \dots, M\}$ 在其 $d$ 维构型空间上
拥有平移自由度。该自由度由一个 $d$ 维 Abel Lie 群描述：

\[
    G_m := (\R^d, +), \quad \g_m := \lie{G}_m \cong \R^d.
\]

$g_m \in G_m$ 的作用方式为：
\[
    T_{g_m}: \R^d \to \R^d, \quad T_{g_m}(x) = x + g_m,
\]
即对构型空间中的点做整体平移。

其 Lie 代数 $\g_m$ 的 Lie 括号恒为零：$[\xi, \eta] = 0$ 对所有 $\xi, \eta \in \g_m$，
故 $G_m$ 为 Abel 群。

\medskip
\noindent
In the SCX framework, each expert $m \in \{1, 2, \dots, M\}$ possesses translational
degrees of freedom on its $d$-dimensional configuration space.
This freedom is described by a $d$-dimensional abelian Lie group:

\vspace{-0.3em}
\[
    G_m := (\R^d, +), \quad \g_m := \lie{G}_m \cong \R^d.
\]

\vspace{-0.3em}
The action of $g_m \in G_m$ is:
$T_{g_m}: \R^d \to \R^d$, $T_{g_m}(x) = x + g_m$,
i.e., a global translation of points in the configuration space.

Its Lie algebra $\g_m$ has identically zero Lie bracket:
$[\xi, \eta] = 0$ for all $\xi, \eta \in \g_m$,
so $G_m$ is abelian.

\subsection{总规范群：专家直积 / Total Gauge Group: Direct Product of Experts}

所有 $M$ 个专家的平移自由度构成总规范群：

\[
    G := \prod_{m=1}^{M} G_m = \underbrace{\R^d \times \R^d \times \cdots \times \R^d}_{M \text{ 个因子}} \cong \R^{M d}.
\]

$G$ 的 Lie 代数为：
\[
    \g := \lie{G} = \bigoplus_{m=1}^{M} \g_m \cong \R^{M d}.
\]

作为 Abel 群，$G$ 中的群运算为逐分量加法：
\[
    g = (g_1, g_2, \dots, g_M) \in G, \quad
    g \cdot h = (g_1 + h_1, \dots, g_M + h_M).
\]

单位元为 $e = (0, 0, \dots, 0)$，逆元为 $g^{-1} = (-g_1, \dots, -g_M)$。

{\bf 关键洞察：} 尽管 $G$ 本身是 Abel 的，但其取值可能依赖于底流形上的位置
（即 $g = g(x)$ 是 $x \in \base$ 的函数），这引入了非平凡的几何结构。

\medskip
\noindent
The translational degrees of freedom of all $M$ experts form the total gauge group:

\vspace{-0.3em}
\[
    G := \prod_{m=1}^{M} G_m = \underbrace{\R^d \times \R^d \times \cdots \times \R^d}_{M \text{ factors}} \cong \R^{M d}.
\]

\vspace{-0.3em}
The Lie algebra of $G$ is:
$\g := \lie{G} = \bigoplus_{m=1}^{M} \g_m \cong \R^{M d}$.

As an abelian group, the group operation in $G$ is componentwise addition:
$g = (g_1, g_2, \dots, g_M) \in G$,
$g \cdot h = (g_1 + h_1, \dots, g_M + h_M)$.

The identity element is $e = (0, 0, \dots, 0)$, the inverse is
$g^{-1} = (-g_1, \dots, -g_M)$.

{\bf Key insight:} Although $G$ itself is abelian, its values may depend on
the base point $x \in \base$ (i.e., $g = g(x)$ is a function on $\base$),
which introduces non-trivial geometric structure.

\subsection{规范群的物理诠释 / Physical Interpretation of the Gauge Group}

每个 $G_m$ 表示专家 $m$ 对其坐标原点的平移自由度。
如果专家 $m$ 的预测输出为 $\tilde{x}_m \in \R^d$，
则其``物理''位置（即规范不变的``真实''构型）为：
\[
    x^{\text{phys}}_m = \tilde{x}_m - g_m,
\]
其中 $g_m$ 是专家 $m$ 的（未知）规范偏移。

不同专家 $i, j$ 之间的规范差定义为：
\[
    \Delta_{ij} := g_i - g_j \in \R^d.
\]
若对所有 $i, j$ 有 $\Delta_{ij} = 0$，则所有专家处于同一规范。

\medskip
\noindent
Each $G_m$ represents the translational freedom of expert $m$ in choosing its coordinate origin.
If expert $m$'s predicted output is $\tilde{x}_m \in \R^d$,
then its ``physical'' position (i.e., gauge-invariant ``true'' configuration) is:
$x^{\text{phys}}_m = \tilde{x}_m - g_m$,
where $g_m$ is the (unknown) gauge offset of expert $m$.

The gauge difference between experts $i, j$ is defined as:
$\Delta_{ij} := g_i - g_j \in \R^d$.
If $\Delta_{ij} = 0$ for all $i, j$, then all experts share the same gauge.

% ══════════════════════════════════════════════════════════════════════
\section{SCX 主纤维丛 / The SCX Principal Bundle}
% ══════════════════════════════════════════════════════════════════════

\subsection{全空间与底流形 / Total Space and Base Manifold}

\begin{definition}[SCX 主丛]
定义 SCX 主纤维丛 $\bun$ 如下：
\begin{itemize}
    \item {\bf 全空间 $\bun$：} $\bun := \R^{Md} \times \mathcal{X}$，
        其中 $\mathcal{X} \subset \R^{K}$ 是参数/输入空间。
        $\bun$ 中的点 $(g_1, \dots, g_M, x)$ 表示所有专家的规范参数与配置参数。
    \item {\bf 结构群：} $G = \prod_{m=1}^M G_m \cong \R^{Md}$，
        右作用为 $R_h(g, x) := (g + h, x)$（逐分量加法）。
    \item {\bf 投影映射：} $\pi: \bun \to \base$ 定义为
        $\pi(g, x) := [g, x]_{\sim}$，
        其中等价关系 $\sim$ 为：$(g, x) \sim (g', x')$ 当且仅当 $x = x'$ 且
        存在 $h \in G$ 使得 $g' = g + h$。
    \item {\bf 底流形 $\base$：} $\base := \bun / G$，即规范轨道空间。
        其元素等价于所有物理上可区分的专家构型。
\end{itemize}
\end{definition}

\noindent
Define the SCX principal fiber bundle $\bun$ as follows:
\begin{itemize}
    \item {\bf Total space $\bun$:} $\bun := \R^{Md} \times \mathcal{X}$,
        where $\mathcal{X} \subset \R^{K}$ is the parameter/input space.
        A point $(g_1, \dots, g_M, x) \in \bun$ represents gauge parameters and
        configuration parameters of all experts.
    \item {\bf Structure group:} $G = \prod_{m=1}^M G_m \cong \R^{Md}$,
        right action $R_h(g, x) := (g + h, x)$ (componentwise addition).
    \item {\bf Projection map:} $\pi: \bun \to \base$ defined as
        $\pi(g, x) := [g, x]_{\sim}$,
        where the equivalence relation $\sim$ is: $(g, x) \sim (g', x')$ iff $x = x'$ and
        there exists $h \in G$ such that $g' = g + h$.
    \item {\bf Base manifold $\base$:} $\base := \bun / G$, the space of gauge orbits.
        Its elements are equivalent to all physically distinguishable expert configurations.
\end{itemize}

\begin{remark}
由于 $G$ 是向量空间 $\R^{Md}$，$\base \cong \mathcal{X} \times (\R^{Md} / \R^{Md}) \cong \mathcal{X}$。
即底流形 $\base$ 实际上同胚于参数空间 $\mathcal{X}$，但在底流形上不同的点可能通过
``扭曲''的纤维结构对应到不同的物理情境。

然而，真正非平凡的结构来自于{\bf 联络}——它决定了不同纤维上的点
如何``水平地''关联。
\end{remark}

\noindent
Since $G$ is the vector space $\R^{Md}$,
$\base \cong \mathcal{X} \times (\R^{Md} / \R^{Md}) \cong \mathcal{X}$.
That is, the base manifold $\base$ is homeomorphic to the parameter space $\mathcal{X}$,
but different points on the base manifold may correspond to different physical
situations through a ``twisted'' fiber structure.

The truly non-trivial structure comes from the {\bf connection} — it determines how
points on different fibers are ``horizontally'' related.

\subsection{规范轨道作为纤维 / Gauge Orbits as Fibers}

\begin{definition}[纤维]
对 $x \in \base$，其上的纤维 $\bun_x := \pi^{-1}(x) \cong G$ 是
所有具有相同物理内容的规范参数的集合——即规范轨道 (gauge orbit)。
\end{definition}

\noindent
For $x \in \base$, the fiber $\bun_x := \pi^{-1}(x) \cong G$ is the set of all
gauge parameters with the same physical content — i.e., the gauge orbit.

\subsection{局部截面与规范选择 / Local Sections as Gauge Choices}

\begin{definition}[SCX 局部截面]
一个局部截面 $s: U \to \bun$（$U \subset \base$ 为开集）给出了一族专家规范的
局部赋值：
\[
    s(x) = (g_1(x), g_2(x), \dots, g_M(x); x) \in \bun.
\]
这在物理上对应于在参数空间区域 $U$ 上``选择了一个规范''。
\end{definition}

\noindent
A local section $s: U \to \bun$ ($U \subset \base$ open) assigns a family of
expert gauges locally:
$s(x) = (g_1(x), g_2(x), \dots, g_M(x); x) \in \bun$.
This physically corresponds to ``choosing a gauge'' on the parameter-space region $U$.

\begin{definition}[全局截面 / Global Section]
若 $s: \base \to \bun$ 在整个底流形上定义，则称 $s$ 为全局截面，
对应于{\bf 全局规范固定 (global gauge-fixing)}。
\end{definition}

\noindent
If $s: \base \to \bun$ is defined on the entire base manifold, then $s$ is called a
global section, corresponding to {\bf global gauge-fixing}.

% ══════════════════════════════════════════════════════════════════════
\section{联络与平行移动 / Connection and Parallel Transport}
% ══════════════════════════════════════════════════════════════════════

\subsection{联络1-形式的构造 / Construction of the Connection 1-Form}

\begin{definition}[SCX 联络]
在 SCX 主丛 $\bun$ 上，Ehresmann 联络 $\conn \in \Omega^1(\bun, \g)$
被定义为对每个专家 $m$ 定义其与``参考规范''的偏离。

由于 $\g \cong \R^{Md}$，我们可以将联络写为分量形式：
\[
    \conn = (\conn_1, \conn_2, \dots, \conn_M), \quad \conn_m \in \Omega^1(\bun, \R^d).
\]

对局部截面 $s$，拉回联络给出规范势 $A^s = s^*\conn \in \Omega^1(U, \R^{Md})$：
\[
    A^s_m(x) = \sum_{\mu=1}^{\dim \base} A^s_{m,\mu}(x) \, \dif x^\mu,
\]
其中 $A^s_m(x) \in \R^d$ 是专家 $m$ 在点 $x$ 处的 $d$ 维向量势。

物理上，$A^s_m(x)$ 度量了专家 $m$ 的规范随参数 $x$ 变化的``速率''。
\end{definition}

\noindent
On the SCX principal bundle $\bun$, the Ehresmann connection
$\conn \in \Omega^1(\bun, \g)$ is defined to measure, for each expert $m$,
its deviation from a ``reference gauge.''

Since $\g \cong \R^{Md}$, we can write the connection in component form:
$\conn = (\conn_1, \conn_2, \dots, \conn_M)$,
$\conn_m \in \Omega^1(\bun, \R^d)$.

For a local section $s$, the pullback connection gives the gauge potential
$A^s = s^*\conn \in \Omega^1(U, \R^{Md})$:
$A^s_m(x) = \sum_{\mu=1}^{\dim \base} A^s_{m,\mu}(x) \, \dif x^\mu$,
where $A^s_m(x) \in \R^d$ is the $d$-dimensional vector potential for expert $m$ at $x$.

Physically, $A^s_m(x)$ measures the ``rate'' at which expert $m$'s gauge varies
with parameters $x$.

\subsection{平行移动与专家对齐 / Parallel Transport and Expert Alignment}

\begin{definition}[沿曲线的平行移动]
设 $\gamma: [0,1] \to \base$ 为光滑曲线，$\gamma(0) = x_0$。
给定纤维 $\bun_{x_0}$ 上的初始点 $p_0 = (g_m^0, x_0)$，
$p_0$ 沿 $\gamma$ 的水平提升 $\tilde{\gamma}: [0,1] \to \bun$ 是满足以下条件的
唯一曲线：
\[
    \pi(\tilde{\gamma}(t)) = \gamma(t), \quad
    \tilde{\gamma}'(t) \in H_{\tilde{\gamma}(t)} \bun, \quad
    \tilde{\gamma}(0) = p_0.
\]

平行移动定义为映射：
\[
    P_\gamma: \bun_{x_0} \to \bun_{x_1}, \quad
    P_\gamma(p_0) := \tilde{\gamma}(1).
\]
\end{definition}

\noindent
Let $\gamma: [0,1] \to \base$ be a smooth curve with $\gamma(0) = x_0$.
Given an initial point $p_0 = (g_m^0, x_0)$ in the fiber $\bun_{x_0}$,
the horizontal lift $\tilde{\gamma}: [0,1] \to \bun$ of $p_0$ along $\gamma$ is the
unique curve satisfying:
$\pi(\tilde{\gamma}(t)) = \gamma(t)$,
$\tilde{\gamma}'(t) \in H_{\tilde{\gamma}(t)} \bun$,
$\tilde{\gamma}(0) = p_0$.

Parallel transport is defined as the map:
$P_\gamma: \bun_{x_0} \to \bun_{x_1}$,
$P_\gamma(p_0) := \tilde{\gamma}(1)$.

\begin{proposition}[平行移动的专家解释]
对于 SCX 主丛，沿 $\gamma$ 从 $x_0$ 到 $x_1$ 的平行移动
将一个专家的规范选择 $g_m(x_0)$ 映射到 $g_m(x_1)$，
其方式要求在参数空间中沿着 $\gamma$ 移动时，
专家之间的{\bf 相对规范保持常数}（相对于联络 $\conn$）。

具体地，水平条件 $\tilde{\gamma}'(t) \in H_{\tilde{\gamma}(t)} \bun$ 等价于：
\[
    \conn_m\left(\frac{\dif g_m}{\dif t}\right) = 0, \quad \forall m = 1, \dots, M.
\]
\end{proposition}

\noindent
For the SCX principal bundle, parallel transport along $\gamma$ from $x_0$ to $x_1$
maps an expert's gauge choice $g_m(x_0)$ to $g_m(x_1)$ in such a way that the
{\bf relative gauge between experts remains constant} (with respect to $\conn$)
when moving along $\gamma$ in parameter space.

Specifically, the horizontality condition
$\tilde{\gamma}'(t) \in H_{\tilde{\gamma}(t)} \bun$ is equivalent to:
$\conn_m(\frac{\dif g_m}{\dif t}) = 0$ for all $m = 1, \dots, M$.

\subsection{联络的分量结构 / Component Structure of the Connection}

由于 $G$ 是 Abel 群，我们可以将联络写为：
\[
    \conn = \sum_{m=1}^{M} \conn_m \otimes e_m,
\]
其中 $\{e_m\}_{m=1}^{M}$ 是 $\g \cong \R^{Md}$ 的标准基，
每个 $\conn_m \in \Omega^1(\bun, \R^d)$ 满足：
\begin{enumerate}[label=(\arabic*)]
    \item $\conn_m(\xi^\sharp) = \xi_m$（$G$-等变的竖直条件）；
    \item $R_h^* \conn_m = \conn_m$（因为 $G$ 是 Abel 的，$\Ad_{g^{-1}} = \id$）。
\end{enumerate}

对局部截面 $s$，我们得到局部 $1$-形式 $A^s_m := s^* \conn_m$。
由于 $G$ 是 Abel 的，规范变换 $g(x) \in G$ 下：
\[
    A^{s'}_m = A^s_m - \dif g_m,
\]
其中 $g = (g_1, \dots, g_M) \in G$。

\medskip
\noindent
Since $G$ is abelian, we can write the connection as:
$\conn = \sum_{m=1}^{M} \conn_m \otimes e_m$,
where $\{e_m\}_{m=1}^{M}$ is the standard basis of $\g \cong \R^{Md}$,
and each $\conn_m \in \Omega^1(\bun, \R^d)$ satisfies:
\begin{enumerate}[label=(\arabic*)]
    \item $\conn_m(\xi^\sharp) = \xi_m$ (vertical $G$-equivariance condition);
    \item $R_h^* \conn_m = \conn_m$ (since $G$ is abelian, $\Ad_{g^{-1}} = \id$).
\end{enumerate}

For a local section $s$, we obtain local $1$-forms $A^s_m := s^* \conn_m$.
Under a gauge transformation $g(x) \in G$, since $G$ is abelian:
$A^{s'}_m = A^s_m - \dif g_m$,
where $g = (g_1, \dots, g_M) \in G$.

% ══════════════════════════════════════════════════════════════════════
\section{曲率与全局规范固定障碍 / Curvature and the Obstruction to Global Gauge-Fixing}
% ══════════════════════════════════════════════════════════════════════

\subsection{曲率的形式 / Form of the Curvature}

\begin{definition}[SCX 曲率]
由于 $G$ 是 Abel 群，SCX 联络 $\conn$ 的曲率 $2$-形式简化为外导数：
\[
    \curv := \dif \conn = \sum_{m=1}^{M} \dif \conn_m \otimes e_m \in \Omega^2(\bun, \g).
\]

在局部截面 $s$ 下，场强为：
\[
    F^s = s^* \curv = \sum_{m=1}^{M} \dif A^s_m \otimes e_m \in \Omega^2(U, \R^{Md}).
\]

对每个专家分量 $m$：
\[
    F^s_m = \dif A^s_m = \sum_{\mu < \nu} F^s_{m,\mu\nu} \, \dif x^\mu \wedge \dif x^\nu,
\]
其中
\[
    F^s_{m,\mu\nu} = \partial_\mu A^s_{m,\nu} - \partial_\nu A^s_{m,\mu}.
\]
\end{definition}

\noindent
Since $G$ is abelian, the curvature $2$-form of the SCX connection $\conn$
simplifies to the exterior derivative:
$\curv := \dif \conn = \sum_{m=1}^{M} \dif \conn_m \otimes e_m \in \Omega^2(\bun, \g)$.

Under a local section $s$, the field strength is:
$F^s = s^* \curv = \sum_{m=1}^{M} \dif A^s_m \otimes e_m \in \Omega^2(U, \R^{Md})$.

For each expert component $m$:
$F^s_m = \dif A^s_m = \sum_{\mu < \nu} F^s_{m,\mu\nu} \, \dif x^\mu \wedge \dif x^\nu$,
where $F^s_{m,\mu\nu} = \partial_\mu A^s_{m,\nu} - \partial_\nu A^s_{m,\mu}$.

\subsection{曲率的规范不变性 / Gauge Invariance of Curvature}

\begin{proposition}[场强的规范不变性]
对 Abel 群 $G$，场强 $F$ 是规范不变量：
\[
    F^{s'}_m = \dif A^{s'}_m = \dif (A^s_m - \dif g_m) = \dif A^s_m - \dif^2 g_m = \dif A^s_m = F^s_m.
\]
即 $F^{s'}_m = F^s_m$。因此，曲率 $\curv$ 是整个丛 $\bun$ 上全局定义的张量场，
而不仅仅是局部截面上的对象。
\end{proposition}

\noindent
For abelian $G$, the field strength $F$ is gauge-invariant:
$F^{s'}_m = \dif A^{s'}_m = \dif (A^s_m - \dif g_m) = \dif A^s_m - \dif^2 g_m = \dif A^s_m = F^s_m$.
That is, $F^{s'}_m = F^s_m$. Thus, the curvature $\curv$ is a globally defined tensor field
on the entire bundle $\bun$, not merely an object on a local section.

\subsection{平坦联络与全局规范固定 / Flat Connection and Global Gauge-Fixing}

\begin{theorem}[全局规范固定定理]
\label{thm:global_gauge}
设 $\bun \to \base$ 为 SCX 主丛，配备联络 $\conn$ 及曲率 $\curv$。
假设底流形 $\base$ 是单连通的。则以下条件等价：
\begin{enumerate}[label=(\roman*)]
    \item $\curv = 0$（联络平坦）；
    \item 存在全局截面 $s: \base \to \bun$，使得 $A^s = s^*\conn = 0$；
    \item 所有专家共享同一个全局规范，即存在一个{\bf 全局规范固定}；
    \item 沿 $\base$ 中任意闭曲线的平行移动是平凡的（和乐群 $\Hol(\conn) = \{e\}$）。
\end{enumerate}
\end{theorem}

\begin{proof}
$(i) \Rightarrow (ii)$：
由于 $\curv = \dif \conn = 0$，联络 $\conn$ 是闭形式。
在单连通底流形 $\base$ 上，Poincaré 引理给出了全局 $0$-形式
$\phi: \bun \to \g$，满足 $\conn = \dif \phi$（这里需做拉回到截面，详见下）。
更标准地使用 Frobenius 定理：
$\curv = 0$ 意味着水平分布 $H\bun = \ker(\conn)$ 是可积的。
对单连通 $\base$，水平提升不依赖于路径，这给出全局水平截面，
即满足 $\conn|_{\text{截面}} = 0$ 的全局截面 $s$。

具体地：固定基点 $x_0 \in \base$，选择纤维 $\bun_{x_0}$ 上一点 $p_0$。
对任意 $x \in \base$，取路径 $\gamma: x_0 \to x$，
定义 $s(x) := P_\gamma(p_0)$（平行移动）。
由于 $\curv = 0$ 且 $\base$ 单连通，$s(x)$ 不依赖于 $\gamma$ 的选择。
按定义，$s$ 是水平的，故 $s^*\conn = 0$。

$(ii) \Rightarrow (iii)$：
截面 $s$ 满足 $A^s_m = s^*\conn_m = 0$ 对所有 $m$ 成立。
这意味着在 $s$ 定义的规范下，所有专家的规范势为零。
由于 $A^s_m$ 度量了专家规范的改变率，
$A^s_m = 0$ 意味着每个专家的规范选择在 $\base$ 上是常数。
因此，存在全局规范使得所有专家共享同一规范原点。

$(iii) \Rightarrow (iv)$：
若存在全局规范固定，则沿任何闭曲线的和乐必然为恒等映射
（因为在全局截面上平行移动是平凡的）。

$(iv) \Rightarrow (i)$：
若和乐群平凡，则平行移动与路径无关。
取无穷小回路，沿其平行移动的偏离由曲率度量：
\[
    P_\gamma(p) - p = -\curv_p(X, Y) \cdot \text{面积} + O(\text{面积}^{3/2}),
\]
其中 $X, Y$ 是回路在 $x$ 处的切向量。
和乐平凡意味着对所有 $X, Y$ 有 $\curv_p(X, Y) = 0$，故 $\curv = 0$。
\end{proof}

\noindent
\medskip
\noindent
{\bf Theorem~\ref{thm:global_gauge} (Global Gauge-Fixing Theorem).}
Let $\bun \to \base$ be the SCX principal bundle equipped with connection $\conn$
and curvature $\curv$. Assume the base manifold $\base$ is simply connected.
Then the following are equivalent:
\begin{enumerate}[label=(\roman*)]
    \item $\curv = 0$ (the connection is flat);
    \item There exists a global section $s: \base \to \bun$ such that $A^s = s^*\conn = 0$;
    \item All experts share the same global gauge, i.e., a {\bf global gauge-fixing} exists;
    \item Parallel transport along any closed curve in $\base$ is trivial
        (holonomy group $\Hol(\conn) = \{e\}$).
\end{enumerate}

\begin{proof}[Proof (English summary)]
$(i) \Rightarrow (ii)$: Since $\curv = \dif \conn = 0$, the connection is closed.
On the simply connected base $\base$, Frobenius' theorem implies the horizontal
distribution $H\bun = \ker(\conn)$ is integrable. The horizontal lift of paths
is path-independent, yielding a global horizontal section $s$ with $s^*\conn = 0$.

$(ii) \Rightarrow (iii)$: $A^s_m = 0$ for all $m$ implies each expert's gauge is
constant across $\base$. Hence a global gauge exists where all experts share
the same origin.

$(iii) \Rightarrow (iv)$: A global gauge-fixing renders parallel transport trivial,
so holonomy must be trivial.

$(iv) \Rightarrow (i)$: Trivial holonomy around infinitesimal loops forces
$\curv_p(X,Y) = 0$ for all $X, Y$, hence $\curv = 0$.
\end{proof}

\subsection{曲率的物理意义：势能面不齐 / Physical Meaning of Curvature: Potential Energy Surface Misalignment}

\begin{theorem}[势能面不齐定理]
\label{thm:misalignment}
若 $\curv \neq 0$，则不存在同时满足所有专家对之间规范对齐的单一规范选择。
换言之，不存在全局赋值 $\{g_m(x)\}_{m=1}^M$ 使得对所有 $i, j$ 和所有的 $x$ 有
$\Delta_{ij}(x) = g_i(x) - g_j(x) = \text{常数}$（不依赖于 $x$）。
等价地，所有专家的势能面（PES）不能通过单一的全局平移而完全对齐。
此即{\bf 势能面不齐 (Potential Energy Surface Misalignment)}。
\end{theorem}

\begin{proof}
假设存在全局赋值 $g(x) = (g_1(x), \dots, g_M(x))$ 使得对所有的 $i, j$ 有
$\dif(g_i - g_j) = 0$。则对于局部截面 $s$ 拉回的联络 $A^s$：

对任意两个专家 $i, j$，规范差 $\Delta_{ij}$ 在 $x$ 上的不变性意味着
存在常数向量 $c_{ij} \in \R^d$ 使得 $g_i(x) - g_j(x) = c_{ij}$。
定义等效的单一规范函数 $\bar{g}: \base \to \R^d$ 为 $\bar{g}(x) := g_1(x)$（取专家1为参考），
则 $g_m(x) = \bar{g}(x) + c_{1m}$。

在此全局规范 $g$（视为截面 $s_g$）下：
\[
    A^{s_g}_m = s_g^* \conn_m.
\]
由于所有 $g_m$ 仅通过一个公共函数 $\bar{g}$ 相差常数，
其外导数满足 $\dif A^{s_g}_m = \dif A^{s_g}_1$ 对所有 $m$ 成立。
这意味着 $F^s_m = \dif A^{s_g}_m$ 对所有 $m$ 相同。

然而，若原始联席的曲率非零，即 $\curv \neq 0$，
则存在某局部截面下 $F^s_i \neq F^s_j$ 对某 $i \neq j$ 成立。
由于 $F$ 是规范不变的（命题5.2），该不等式在任何规范下都成立，
与上述全局对齐要求的 $F^s_i = F^s_j$ 矛盾。

因此，$\curv \neq 0$ 蕴含不存在全局规范对齐。
\end{proof}

\noindent
\medskip
\noindent
{\bf Theorem~\ref{thm:misalignment} (Potential Energy Surface Misalignment Theorem).}
If $\curv \neq 0$, then no single gauge choice can simultaneously align all pairwise
expert comparisons. Equivalently, there does not exist a global assignment
$\{g_m(x)\}_{m=1}^M$ such that
$\Delta_{ij}(x) = g_i(x) - g_j(x) = \text{constant}$ (independent of $x$)
for all $i,j$ and all $x$. This means the potential energy surfaces (PES) of all
experts cannot be simultaneously aligned by a single global translation.
This is the phenomenon of {\bf 势能面不齐 (Potential Energy Surface Misalignment)}.

\begin{proof}[Proof (English)]
Assume a global assignment $g(x) = (g_1(x), \dots, g_M(x))$ exists with
$\dif(g_i - g_j) = 0$ for all $i,j$. Then there exist constants $c_{ij}$ such that
$g_i(x) - g_j(x) = c_{ij}$. Choosing expert 1 as reference,
$g_m(x) = \bar{g}(x) + c_{1m}$ for all $m$, where $\bar{g}(x) := g_1(x)$.

Under this global gauge $g$ (viewed as section $s_g$),
$A^{s_g}_m = s_g^* \conn_m$. Since all $g_m$ differ only by constants from $\bar{g}$,
$\dif A^{s_g}_m = \dif A^{s_g}_1$ for all $m$, implying
$F^{s_g}_i = F^{s_g}_j$ for all $i,j$.

But if the original connection has non-zero curvature, there exists some local section
under which $F^s_i \neq F^s_j$ for some $i \neq j$. Since $F$ is gauge-invariant
(Proposition~5.2), this inequality holds in any gauge, contradicting the global
alignment requirement $F_i = F_j$.

Hence $\curv \neq 0$ implies no global gauge alignment exists.
\end{proof}

\begin{remark}[几何解释]
在几何上，$\curv \neq 0$ 意味着主丛 $\bun$ 是``弯曲''的。
即使底流形 $\base$ 是平凡的（同胚于 $\R^K$），
其上的联络可以具有非零曲率，从而引入非平凡的和乐。
这导致了``平行移动依赖于路径''的现实，使得不同专家
沿不同路径在参数空间中移动时，其规范的比较产生内在矛盾。

这正是 Yang-Mills 理论中非零场强对应于非平凡物理（如电磁场、胶子场）的类似现象。
在 SCX 中，$\curv \neq 0$ 表示多专家系统具有{\bf 内在的、不可消除的规范不一致性}。
\end{remark}

\noindent
Geometrically, $\curv \neq 0$ means the principal bundle $\bun$ is ``curved.''
Even if the base manifold $\base$ is trivial (homeomorphic to $\R^K$),
the connection on it can have non-zero curvature, introducing non-trivial holonomy.
This leads to the reality that parallel transport depends on the path,
creating an intrinsic contradiction in comparing different experts' gauges
when moving along different paths in parameter space.

This is precisely analogous to how non-zero field strength in Yang-Mills theory
corresponds to non-trivial physics (electromagnetic fields, gluon fields).
In SCX, $\curv \neq 0$ indicates that the multi-expert system possesses
{\bf intrinsic, ineliminable gauge inconsistency}.

% ══════════════════════════════════════════════════════════════════════
\section{曲率作为拓扑障碍类 / Curvature as a Topological Obstruction Class}
% ══════════════════════════════════════════════════════════════════════

\subsection{和乐与曲率 / Holonomy and Curvature}

\begin{theorem}[和乐-曲率关系]
对 SCX 主丛（$G$ Abel）上的联络 $\conn$，围绕基点 $x_0 \in \base$ 上
无穷小平行四边形（边由切向量 $X, Y \in T_{x_0}\base$ 生成）的和乐为：
\[
    \Hol_{\gamma_{X,Y}}(p_0) = p_0 \cdot \exp\left(-\curv_{p_0}(\tilde{X}, \tilde{Y}) \, \varepsilon^2 + O(\varepsilon^3)\right),
\]
其中 $\tilde{X}, \tilde{Y}$ 是 $X, Y$ 的提升，$\varepsilon$ 是回路尺度。
对 Abel 群：
\[
    \text{和乐平移} = -\curv_{p_0}(\tilde{X}, \tilde{Y}) \, \varepsilon^2.
\]
\end{theorem}

\noindent
For a connection $\conn$ on the SCX principal bundle ($G$ abelian),
the holonomy around an infinitesimal parallelogram based at $x_0 \in \base$
(spanned by tangent vectors $X, Y \in T_{x_0}\base$) is:
$\Hol_{\gamma_{X,Y}}(p_0) = p_0 \cdot \exp(-\curv_{p_0}(\tilde{X}, \tilde{Y}) \, \varepsilon^2 + O(\varepsilon^3))$,
where $\tilde{X}, \tilde{Y}$ are lifts of $X, Y$, and $\varepsilon$ is the loop scale.
For abelian $G$:
$\text{holonomy translation} = -\curv_{p_0}(\tilde{X}, \tilde{Y}) \, \varepsilon^2$.

\subsection{曲率的 de Rham 上同调类 / The de Rham Cohomology Class of Curvature}

\begin{proposition}[曲率作为特征类]
对于 Abel 结构群 $G$，曲率 $\curv$ 的每个分量 $\curv_m \in \Omega^2(\bun, \R^d)$
是闭形式（$\dif \curv_m = 0$，由 Bianchi 恒等式）。
在底流形 $\base$ 上拉回后，$F^s_m = s^*\curv_m$ 定义了 de Rham 上同调类
\[
    [F_m] \in H^2_{\mathrm{dR}}(\base; \R^d).
\]
该上同调类是规范不变的，且为非零当且仅当主丛是非平凡的和/或联络具有非零曲率。
\end{proposition}

\noindent
For abelian structure group $G$, each component $\curv_m \in \Omega^2(\bun, \R^d)$
of the curvature is closed ($\dif \curv_m = 0$, by Bianchi identity).
After pullback to the base manifold $\base$, $F^s_m = s^*\curv_m$ defines de Rham
cohomology classes $[F_m] \in H^2_{\mathrm{dR}}(\base; \R^d)$.
These classes are gauge-invariant and non-zero precisely when the principal bundle
is non-trivial and/or the connection has non-zero curvature.

\begin{corollary}[上同调障碍]
若 $H^2_{\mathrm{dR}}(\base; \R^d) = 0$（例如 $\base$ 可缩），则任何联络的曲率
均可写为 $\curv = \dif \conn$ 且 $\conn$ 可以全局定义，
但这不保证 $\curv = 0$。
全局规范固定要求更强的条件 $\curv = 0$（而非仅 $\curv$ 正合）。
\end{corollary}

\noindent
If $H^2_{\mathrm{dR}}(\base; \R^d) = 0$ (e.g., $\base$ is contractible), then the
curvature of any connection can be written as $\curv = \dif \conn$ with
$\conn$ globally defined, but this does not guarantee $\curv = 0$.
Global gauge-fixing requires the stronger condition $\curv = 0$
(rather than merely $\curv$ being exact).

\subsection{离散情形的曲率：环绕和乐 / Discrete Curvature: Loop Holonomy}

在实际的 SCX 计算中，参数空间通常是离散的（有限的测试构型集）。
此时曲率通过有限回路上的和乐来度量。

\begin{definition}[离散曲率算子]
设 $\mathcal{G} = (\mathcal{V}, \mathcal{E})$ 为底流形 $\base$ 的离散化（图），
其中顶点 $v \in \mathcal{V}$ 对应参数构型，边 $e \in \mathcal{E}$ 对应参数之间的跃迁。
对每条边 $e = (v \to w)$，平行移动定义了群元 $P_e \in G$。

对图中任意有向回路 $\gamma = (v_0 \to v_1 \to \cdots \to v_k = v_0)$，
离散曲率（和乐）为：
\[
    \curv_\gamma := P_{v_{k-1} \to v_0} \circ \cdots \circ P_{v_0 \to v_1} \in G.
\]
对 Abel 群 $G$：
\[
    \curv_\gamma = \sum_{i=0}^{k-1} \Delta g_{(v_i \to v_{i+1})},
\]
其中 $\Delta g_e \in \R^{Md}$ 是边 $e$ 上的规范平移。
\end{definition}

\noindent
In practical SCX computation, the parameter space is typically discrete
(a finite set of test configurations). Curvature is then measured by holonomy
around finite loops.

Let $\mathcal{G} = (\mathcal{V}, \mathcal{E})$ be a discretization (graph) of the base
manifold $\base$, where vertices $v \in \mathcal{V}$ correspond to parameter configurations
and edges $e \in \mathcal{E}$ correspond to transitions between parameters.
For each edge $e = (v \to w)$, parallel transport defines a group element $P_e \in G$.

For any directed loop $\gamma = (v_0 \to v_1 \to \cdots \to v_k = v_0)$ in the graph,
the discrete curvature (holonomy) is:
$\curv_\gamma := P_{v_{k-1} \to v_0} \circ \cdots \circ P_{v_0 \to v_1} \in G$.
For abelian $G$:
$\curv_\gamma = \sum_{i=0}^{k-1} \Delta g_{(v_i \to v_{i+1})}$,
where $\Delta g_e \in \R^{Md}$ is the gauge translation on edge $e$.

\begin{proposition}[离散平坦性判别]
离散联络是平坦的（$\curv_\gamma = 0$ 对所有回路 $\gamma$）当且仅当存在
顶点的全局赋值 $g: \mathcal{V} \to G$ 使得对所有边 $e = (v \to w)$：
\[
    P_e = g(w) - g(v) \quad (\text{在 Abel 群中}).
\]
此即离散版本的 Poincaré 引理：闭 $1$-形式（在图上无边）是正合的。
\end{proposition}

\noindent
A discrete connection is flat ($\curv_\gamma = 0$ for all loops $\gamma$)
iff there exists a global vertex assignment $g: \mathcal{V} \to G$ such that
for all edges $e = (v \to w)$: $P_e = g(w) - g(v)$ (in the abelian group).
This is the discrete version of the Poincaré lemma:
a closed $1$-form (curl-free on the graph) is exact.

% ══════════════════════════════════════════════════════════════════════
\section{Cercis 分数：规范不变几何可观测量 / The Cercis Score: A Gauge-Invariant Geometric Observable}
% ══════════════════════════════════════════════════════════════════════

\subsection{动机：Yang-Mills 理论中的规范不变量 / Motivation: Gauge Invariants in Yang-Mills}

在 Yang-Mills 理论中，局部规范不变量由场强张量 $F_{\mu\nu}$ 的缩并构成，
如拉格朗日密度 $\Tr(F_{\mu\nu}F^{\mu\nu})$。
这些量是物理上可观测的，因为它们不依赖于规范选择。

在 SCX 框架中，我们需要类似的规范不变量来客观度量
多专家系统的一致性，而不依赖于任何特定的规范固定方案。

\medskip
\noindent
In Yang-Mills theory, local gauge invariants are constructed from contractions of
the field strength tensor $F_{\mu\nu}$, such as the Lagrangian density
$\Tr(F_{\mu\nu}F^{\mu\nu})$. These quantities are physically observable because they
are independent of gauge choices.

In the SCX framework, we need analogous gauge-invariant quantities to objectively
measure the consistency of a multi-expert system, independent of any particular
gauge-fixing scheme.

\subsection{Cercis 分数的定义 / Definition of the Cercis Score}

\begin{definition}[Cercis 分数]
SCX 的 Cercis 分数 $\cercis$ 定义为曲率 $2$-形式的 $L^2$ 范数（在底流形上的积分）：
\[
    \cercis := \int_{\base} \|F\|^2 \, \vol_{\base},
\]
其中
\[
    \|F\|^2 := \sum_{m=1}^{M} \sum_{\mu < \nu} \|F_{m,\mu\nu}\|^2_{\R^d},
\]
$\vol_{\base}$ 是 $\base$ 上的体积形式（或离散情形的计数测度），
且 $F_{m,\mu\nu} = \partial_\mu A_{m,\nu} - \partial_\nu A_{m,\mu}$ 为某（任意）规范下的场强分量。

在离散情形（有限测试构型集 $\mathcal{D} = \{x^1, \dots, x^N\}$）：
\[
    \cercis := \frac{1}{|\mathcal{L}|} \sum_{\gamma \in \mathcal{L}} \|\curv_\gamma\|^2,
\]
其中 $\mathcal{L}$ 是离散化底流形中所有基本回路的集合，$\curv_\gamma$ 是沿回路 $\gamma$ 的离散曲率。
\end{definition}

\noindent
The SCX Cercis Score $\cercis$ is defined as the $L^2$ norm of the curvature $2$-form
(integrated over the base manifold):
$\cercis := \int_{\base} \|F\|^2 \, \vol_{\base}$,
where
$\|F\|^2 := \sum_{m=1}^{M} \sum_{\mu < \nu} \|F_{m,\mu\nu}\|^2_{\R^d}$,
$\vol_{\base}$ is the volume form on $\base$ (or counting measure in the discrete case),
and $F_{m,\mu\nu} = \partial_\mu A_{m,\nu} - \partial_\nu A_{m,\mu}$ are the field strength
components in some (arbitrary) gauge.

In the discrete case (finite test configuration set $\mathcal{D} = \{x^1, \dots, x^N\}$):
$\cercis := \frac{1}{|\mathcal{L}|} \sum_{\gamma \in \mathcal{L}} \|\curv_\gamma\|^2$,
where $\mathcal{L}$ is the set of all elementary loops in the discretized base manifold,
and $\curv_\gamma$ is the discrete curvature along loop $\gamma$.

\subsection{Cercis 的规范不变性 / Gauge Invariance of Cercis}

\begin{theorem}[Cercis 的规范不变性]
Cercis 分数 $\cercis$ 是规范不变的：
对任意规范变换 $g: \base \to G$，有 $\cercis' = \cercis$。
\end{theorem}

\begin{proof}
由命题 5.2，$F^{s'}_m = F^s_m$ 对任意两个局部截面 $s, s'$ 成立。
即场强张量 $F$ 是规范不变的。
由于 $\cercis$ 的定义仅依赖于 $F$（通过其 $L^2$ 范数的积分），
$\cercis$ 是规范不变的。$\square$
\end{proof}

\noindent
{\bf Theorem (Gauge Invariance of Cercis).}
The Cercis Score $\cercis$ is gauge-invariant:
for any gauge transformation $g: \base \to G$, $\cercis' = \cercis$.

{\bf Proof.} By Proposition~5.2, $F^{s'}_m = F^s_m$ for any two local sections $s, s'$.
That is, the field strength tensor $F$ is gauge-invariant.
Since $\cercis$ is defined solely in terms of $F$ (via the integral of its $L^2$ norm),
$\cercis$ is gauge-invariant. $\square$

\subsection{Cercis 的几何解释 / Geometric Interpretation of Cercis}

\begin{proposition}[Cercis 作为弯曲度量]
$\cercis = 0$ 当且仅当 $\curv = 0$（在 $\base$ 上几乎处处成立）。
因此：
\begin{itemize}
    \item $\cercis = 0$：多专家系统是{\bf 规范一致的 (gauge-consistent)}——存在全局规范固定；
    \item $\cercis > 0$：系统具有{\bf 不可消除的规范不一致性}——$\cercis$ 度量了不一致的程度。
\end{itemize}
\end{proposition}

\noindent
$\cercis = 0$ if and only if $\curv = 0$ (almost everywhere on $\base$).
Thus:
\begin{itemize}
    \item $\cercis = 0$: the multi-expert system is {\bf gauge-consistent} — global gauge-fixing exists;
    \item $\cercis > 0$: the system possesses {\bf ineliminable gauge inconsistency} —
        $\cercis$ measures the degree of inconsistency.
\end{itemize}

\begin{remark}[与 Yang-Mills 作用量的类比]
Cercis $\cercis = \int \|F\|^2 \vol$ 在形式上等价于 Yang-Mills 理论中的作用量
$S_{\text{YM}} = \int \Tr(F \wedge \star F)$（差一个 Hodge 星算子）。
正如 Yang-Mills 作用量的极小值对应于经典运动方程的解（Yang-Mills 瞬子等），
Cercis 的极小值对应于多专家系统的最优规范调整方案。
\end{remark}

\noindent
The Cercis score $\cercis = \int \|F\|^2 \vol$ is formally analogous to the Yang-Mills
action $S_{\text{YM}} = \int \Tr(F \wedge \star F)$ (up to a Hodge star operator).
Just as the minima of the Yang-Mills action correspond to solutions of the classical
equations of motion (Yang-Mills instantons, etc.), the minima of Cercis correspond
to optimal gauge adjustment schemes for the multi-expert system.

% ══════════════════════════════════════════════════════════════════════
\section{MILP 规范固定：离散 Coulomb 规范 / MILP Gauge-Fixing: Discrete Coulomb Gauge}
% ══════════════════════════════════════════════════════════════════════

\subsection{连续 Coulomb 规范 / The Continuous Coulomb Gauge}

在经典的 Yang-Mills 理论中，Coulomb 规范（也称为辐射规范或横场规范）由条件
\[
    \nabla \cdot \mathbf{A} = 0
\]
定义，即规范势 $A$ 的散度为零。
在几何语言中，这等价于要求联络 $1$-形式 $\conn$ 满足一个水平性条件：
其相对于某个 Riemann 度量的共轭导数为零。

更确切地说，Coulomb 规范是截面 $s$ 的选择（即 $A = s^*\conn$），
使得 $A$ 与底流形上某个固定 Riemann 度量的 Levi-Civita 联络正交。

\medskip
\noindent
In classical Yang-Mills theory, the Coulomb gauge (also called radiation gauge or
transverse gauge) is defined by the condition $\nabla \cdot \mathbf{A} = 0$,
i.e., the divergence of the gauge potential $A$ is zero.
In geometric language, this is equivalent to requiring the connection $1$-form $\conn$
to satisfy a horizontality condition: its co-derivative with respect to some Riemannian
metric vanishes.

More precisely, the Coulomb gauge is a choice of section $s$ (i.e., $A = s^*\conn$)
such that $A$ is orthogonal to the Levi-Civita connection of some fixed Riemannian metric
on the base manifold.

\subsection{Coulomb 规范的变分特征 / Variational Characterization of Coulomb Gauge}

\begin{proposition}[Coulomb 规范的变分原理]
设 $\base$ 配备 Riemann 度量，联络 $1$-形式 $A$ 满足 Coulomb 条件
$\delta A = \nabla^\mu A_\mu = 0$ 当且仅当 $A$ 在所有规范变换下极小化
$L^2$ 范数：
\[
    \|A\|_{L^2}^2 := \int_{\base} \|A\|^2 \, \vol_{\base}.
\]
即 Coulomb 规范是以下变分问题的解：
\[
    \min_{g: \base \to G} \int_{\base} \|A - \dif g\|^2 \, \vol_{\base}.
\]
\end{proposition}

\begin{proof}
考虑泛函 $\mathcal{F}[g] := \int_{\base} \|A - \dif g\|^2 \, \vol_{\base}$。
取变分 $g \mapsto g + \varepsilon \eta$（$\eta: \base \to \g$）：
\[
    \frac{\dif}{\dif \varepsilon}\Big|_{\varepsilon=0} \mathcal{F}[g + \varepsilon \eta]
    = -2 \int_{\base} \langle A - \dif g, \dif \eta \rangle \, \vol_{\base}
    = 2 \int_{\base} \langle \delta(A - \dif g), \eta \rangle \, \vol_{\base}.
\]
其中 $\delta$ 是余微分算子（$\dif$ 的形式伴随，$\delta = \star \dif \star$）。
极值条件 $\delta \mathcal{F} = 0$ 对所有 $\eta$ 成立当且仅当
$\delta(A - \dif g) = 0$，即 $\delta A - \delta \dif g = 0$。

若取规范 $g$ 使得 $\dif g = A$（即选择该规范），则条件简化为 $\delta \dif g = 0$。
使用 Hodge 分解，每个 $1$-形式有唯一分解 $A = \dif \phi + \delta \psi + h$（调和部分）。
Coulomb 规范对应于选取 $\dif \phi = 0$，即去除纵向部分。
\end{proof}

\noindent
Let $\base$ be equipped with a Riemannian metric. The connection $1$-form $A$ satisfies
the Coulomb condition $\delta A = \nabla^\mu A_\mu = 0$ iff $A$ minimizes the $L^2$ norm
among all gauge transformations:
$\|A\|_{L^2}^2 := \int_{\base} \|A\|^2 \, \vol_{\base}$.
That is, the Coulomb gauge solves:
$\min_{g: \base \to G} \int_{\base} \|A - \dif g\|^2 \, \vol_{\base}$.

{\bf Proof (sketch).} Consider the functional
$\mathcal{F}[g] := \int_{\base} \|A - \dif g\|^2 \, \vol_{\base}$.
The first variation yields $\delta(A - \dif g) = 0$ as the stationarity condition,
which is the Coulomb condition.

\subsection{MILP 作为离散 Coulomb 规范 / MILP as Discrete Coulomb Gauge}

在 SCX 的离散实现中，参数空间由有限集 $\{x^k\}_{k=1}^N$ 表示（例如 DFT 计算中的不同构型）。
规范变换 $g_m(x^k) \in \R^d$ 对每个专家 $m$ 和每个构型 $k$ 定义了一个平移。

MILP (Mixed-Integer Linear Programming) 规范固定方案在离散设定中实现了 Coulomb 规范的类比：

\begin{definition}[MILP 规范固定]
SCX 的 MILP 规范固定问题定义为：
\[
\begin{aligned}
    \min_{\{g_m^k\}} \quad & \sum_{m=1}^{M} \sum_{k=1}^{N} \sum_{\mu=1}^{\dim \base}
        \|A_{m,\mu}(x^k) - (\Delta_\mu g_m)(x^k)\|^2 \\
    \text{s.t.} \quad & \sum_{m=1}^{M} g_m^k = 0, \quad \forall k = 1, \dots, N,
\end{aligned}
\]
其中 $\Delta_\mu$ 是参数空间中的离散差分算子（沿第 $\mu$ 个方向的有限差分）。

约束条件 $\sum_{m=1}^{M} g_m^k = 0$ 是连续 Coulomb 条件 $\delta A = 0$ 的离散类比，
它固定了``质心''规范自由度。
\end{definition}

\noindent
In the discrete implementation of SCX, the parameter space is represented by a finite
set $\{x^k\}_{k=1}^N$ (e.g., different configurations in DFT calculations).
The gauge transformation $g_m(x^k) \in \R^d$ defines a translation for each expert $m$
at each configuration $k$.

The MILP (Mixed-Integer Linear Programming) gauge-fixing scheme realizes the analog of
the Coulomb gauge in the discrete setting:

The SCX MILP gauge-fixing problem is defined as:
\[
\begin{aligned}
    \min_{\{g_m^k\}} \quad & \sum_{m=1}^{M} \sum_{k=1}^{N} \sum_{\mu=1}^{\dim \base}
        \|A_{m,\mu}(x^k) - (\Delta_\mu g_m)(x^k)\|^2 \\
    \text{s.t.} \quad & \sum_{m=1}^{M} g_m^k = 0, \quad \forall k = 1, \dots, N,
\end{aligned}
\]
where $\Delta_\mu$ is the discrete difference operator in parameter space
(finite difference along the $\mu$-th direction).
The constraint $\sum_{m=1}^{M} g_m^k = 0$ is the discrete analog of the continuous
Coulomb condition $\delta A = 0$, fixing the ``center-of-mass'' gauge degree of freedom.

\begin{theorem}[MILP-Coulomb 对应定理]
\label{thm:milp_coulomb}
MILP 规范固定方案（含约束 $\sum_m g_m^k = 0$）是连续 Coulomb 规范
在以下意义上的离散精确类比：
\begin{enumerate}[label=(\roman*)]
    \item MILP 极小化了离散联络的 $L^2$ 范数（等价于连续 Coulomb 规范的变分原理）；
    \item 约束 $\sum_m g_m^k = 0$ 对应于连续散度条件 $\nabla^\mu A_\mu = 0$ 的离散形式，
        其中散度算子在离散图中被替换为关于专家的求和；
    \item 在细网格极限下，MILP 的解收敛于连续 Coulomb 规范固定问题的解。
\end{enumerate}
\end{theorem}

\begin{proof}[证明概要]
(i) MILP 的目标函数是连续泛函 $\int \|A - \dif g\|^2 \vol$ 的离散化
（积分 → 求和，导数 → 有限差分）。
极小化该离散泛函等价于极小化离散 $L^2$ 范数，即连续 Coulomb 规范的变分特征。

(ii) 在连续情形，Coulomb 条件 $\nabla^\mu A_\mu = 0$ 使用 Hodge 分解去除了
$A$ 的纵向部分（$\dif \phi$ 项）。
在离散情形，$A$ 的纵向部分对应于所有专家共享的``公共平移''——即形如
$g_m^k = \bar{g}^k$ 对所有 $m$ 的规范变换。
约束 $\sum_m g_m^k = 0$ 精确地移除了这个公共模式，
等价于要求规范变换的``离散散度''为零。

(iii) 当参数空间的离散化无限细化时（$N \to \infty$，网格间距 $\to 0$），
离散 MILP 目标函数逐点收敛于连续 $\|A - \dif g\|^2$ 的积分，
且约束的极限给出 $\int (\sum_m g_m) = 0$，与 Coulomb 规范相容。
\end{proof}

\noindent
\medskip
\noindent
{\bf Theorem~\ref{thm:milp_coulomb} (MILP-Coulomb Correspondence).}
The MILP gauge-fixing scheme (with constraint $\sum_m g_m^k = 0$) is the discrete
exact analog of the continuous Coulomb gauge in the following sense:
\begin{enumerate}[label=(\roman*)]
    \item MILP minimizes the $L^2$ norm of the discrete connection (equivalent to the
        variational principle of the continuous Coulomb gauge);
    \item The constraint $\sum_m g_m^k = 0$ corresponds to the discrete form of the
        continuous divergence condition $\nabla^\mu A_\mu = 0$, where the divergence
        operator is replaced by summation over experts on the discrete graph;
    \item In the fine-grid limit, the MILP solution converges to the solution of the
        continuous Coulomb gauge-fixing problem.
\end{enumerate}

\begin{proof}[Proof (English sketch)]
(i) The MILP objective function is a discretization of the continuous functional
$\int \|A - \dif g\|^2 \vol$ (integral $\to$ sum, derivative $\to$ finite difference).
Minimizing this discrete functional is equivalent to minimizing the discrete $L^2$ norm,
i.e., the variational characterization of the continuous Coulomb gauge.

(ii) In the continuous case, the Coulomb condition $\nabla^\mu A_\mu = 0$ uses the Hodge
decomposition to remove the longitudinal part ($\dif \phi$ term) of $A$.
In the discrete case, the longitudinal part of $A$ corresponds to the ``common translation''
shared by all experts — i.e., gauge transformations of the form $g_m^k = \bar{g}^k$ for all $m$.
The constraint $\sum_m g_m^k = 0$ precisely removes this common mode,
equivalent to requiring the ``discrete divergence'' of the gauge transformation to be zero.

(iii) As the discretization of the parameter space is infinitely refined
($N \to \infty$, grid spacing $\to 0$), the discrete MILP objective converges
pointwise to the integral of the continuous $\|A - \dif g\|^2$,
and the limit of the constraint yields $\int (\sum_m g_m) = 0$, compatible with
the Coulomb gauge.
\end{proof}

\subsection{MILP 的几何解释 / Geometric Interpretation of MILP}

\begin{remark}[规范固定作为截面选择]
在几何上，MILP 规范固定对应于选择主丛 $\bun$ 的一个特定截面 $s_{\text{MILP}}$，
使得拉回联络 $A^{\text{MILP}} = s_{\text{MILP}}^*\conn$ 满足：
\begin{enumerate}[label=(\arabic*)]
    \item $A^{\text{MILP}}$ 的离散 $L^2$ 范数被极小化；
    \item 截面 $s_{\text{MILP}}$ 在纤维 $\bun_{x^k}$ 上的限制满足
        ``质心''条件 $\sum_m g_m^k = 0$。
\end{enumerate}
这精确地定义了纤维丛上的一个（离散）``水平''截面选择，
类似于连续 Coulomb 规范中要求 $A$ 与纵向方向正交。
\end{remark}

\noindent
Geometrically, MILP gauge-fixing corresponds to choosing a specific section
$s_{\text{MILP}}$ of the principal bundle $\bun$ such that the pullback connection
$A^{\text{MILP}} = s_{\text{MILP}}^*\conn$ satisfies:
\begin{enumerate}[label=(\arabic*)]
    \item The discrete $L^2$ norm of $A^{\text{MILP}}$ is minimized;
    \item The restriction of the section $s_{\text{MILP}}$ to fibers $\bun_{x^k}$
        satisfies the ``center-of-mass'' condition $\sum_m g_m^k = 0$.
\end{enumerate}
This precisely defines a (discrete) ``horizontal'' section choice on the fiber bundle,
analogous to requiring $A$ to be orthogonal to the longitudinal direction in the
continuous Coulomb gauge.

% ══════════════════════════════════════════════════════════════════════
\section{变分定理：$\sum g_m = 0$ 最小化联络范数 / Variational Theorem: $\sum g_m = 0$ Minimizes Connection Norm}
% ══════════════════════════════════════════════════════════════════════

\subsection{问题陈述 / Problem Statement}

本节我们证明一个核心变分定理：
在 MILP 中使用的约束条件 $\sum_{m=1}^{M} g_m = 0$ 不是任意的，
而是最小化联络 $\conn$ 的 $L^2$ 范数的自然结果。

\medskip
\noindent
We prove a central variational theorem:
the constraint $\sum_{m=1}^{M} g_m = 0$ used in MILP is not arbitrary,
but is the natural consequence of minimizing the $L^2$ norm of the connection $\conn$.

\subsection{连续版本 / Continuous Version}

\begin{theorem}[联络范数极小化定理 — 连续版本]
\label{thm:min_norm}
设 $\base$ 为紧致 Riemann 流形，$A = s^*\conn \in \Omega^1(\base, \R^{Md})$
为某截面 $s$ 下的规范势。
考虑以下变分问题：
\[
    \min_{g \in C^\infty(\base, G)} \mathcal{I}[g], \quad
    \mathcal{I}[g] := \int_{\base} \|A - \dif g\|^2 \, \vol_{\base},
\]
其中 $\|A - \dif g\|^2 = \sum_{m=1}^M \sum_{\mu=1}^{\dim \base} \|A_{m,\mu} - \partial_\mu g_m\|^2_{\R^d}$。

则该变分问题的 Euler-Lagrange 方程为：
\[
    \delta(A - \dif g) = 0,
\]
其中 $\delta$ 是余微分算子。
在适当边界条件下，其解满足{\bf 库仑条件}：
\[
    \sum_{m=1}^{M} \nabla^\mu (A_{m,\mu} - \partial_\mu g_m) = 0.
\]

特别地，若我们进一步要求解的唯一性（规范固定），自然附加条件
$\sum_{m=1}^{M} g_m = 0$ 以消除``零模''（所有专家整体平移的自由度），
该条件是使得变分问题有唯一解的最小约束。
\end{theorem}

\begin{proof}
变分计算：
\[
\begin{aligned}
    \frac{\dif}{\dif \varepsilon}\Big|_{\varepsilon=0} \mathcal{I}[g + \varepsilon \eta]
    &= \frac{\dif}{\dif \varepsilon}\Big|_{\varepsilon=0}
        \int_{\base} \sum_{m,\mu} \|A_{m,\mu} - \partial_\mu(g_m + \varepsilon \eta_m)\|^2 \, \vol_{\base} \\
    &= -2 \int_{\base} \sum_{m,\mu} (A_{m,\mu} - \partial_\mu g_m) \cdot \partial_\mu \eta_m \, \vol_{\base} \\
    &= 2 \int_{\base} \sum_{m,\mu} \partial_\mu (A_{m,\mu} - \partial_\mu g_m) \cdot \eta_m \, \vol_{\base} \quad \text{(分部积分)}.
\end{aligned}
\]

要求 $\frac{\dif}{\dif \varepsilon}\big|_0 \mathcal{I}[g + \varepsilon \eta] = 0$
对所有变分 $\eta$ 成立，得到 Euler-Lagrange 方程：
\[
    \partial_\mu (A_{m,\mu} - \partial_\mu g_m) = 0, \quad \forall m = 1, \dots, M.
\]

这正是连续 Coulomb 规范条件 $\nabla^\mu (A_\mu - \partial_\mu g) = 0$。

然而，上述方程对 $g$ 并未完全确定：若 $g$ 是一个解，则
$g + c$（其中 $c \in \R^{Md}$ 为常数）也是解。
此即规范固定中残留的``整体平移''自由度。
施加条件 $\int_{\base} \sum_m g_m \, \vol_{\base} = 0$（或其强化形式
$\sum_m g_m(x) = 0$ 对所有 $x$）可以消除该零模。
\end{proof}

\noindent
\medskip
\noindent
{\bf Theorem~\ref{thm:min_norm} (Connection Norm Minimization — Continuous Version).}
Let $\base$ be a compact Riemannian manifold and $A = s^*\conn \in \Omega^1(\base, \R^{Md})$
the gauge potential under some section $s$.
Consider the variational problem:
$\min_{g \in C^\infty(\base, G)} \mathcal{I}[g]$,
$\mathcal{I}[g] := \int_{\base} \|A - \dif g\|^2 \, \vol_{\base}$,
where $\|A - \dif g\|^2 = \sum_{m=1}^M \sum_{\mu=1}^{\dim \base} \|A_{m,\mu} - \partial_\mu g_m\|^2_{\R^d}$.

The Euler-Lagrange equation of this variational problem is $\delta(A - \dif g) = 0$,
where $\delta$ is the co-differential operator. Under appropriate boundary conditions,
its solutions satisfy the {\bf Coulomb condition}:
$\sum_{m=1}^{M} \nabla^\mu (A_{m,\mu} - \partial_\mu g_m) = 0$.

In particular, to ensure uniqueness of solution (gauge-fixing), one naturally imposes
the additional condition $\sum_{m=1}^{M} g_m = 0$ to eliminate the ``zero mode''
(the overall translation degree of freedom of all experts). This condition is the
minimal constraint that makes the variational problem have a unique solution.

\subsection{离散版本 / Discrete Version}

\begin{theorem}[联络范数极小化定理 — 离散版本]
\label{thm:min_norm_discrete}
设 $\mathcal{G} = (\mathcal{V}, \mathcal{E})$ 为底流形 $\base$ 的离散化图。
定义离散联络势 $A: \mathcal{E} \to \R^{Md}$ 为边上的赋值。
考虑离散变分问题：
\[
    \min_{g: \mathcal{V} \to \R^{Md}} \sum_{e=(v \to w) \in \mathcal{E}}
        \|A_e - (g(w) - g(v))\|^2.
\]

该问题的极值条件为（对每个顶点 $v$ 和专家 $m$）：
\[
    \sum_{w \sim v} (A_{m,(v \to w)} - (g_m(w) - g_m(v))) = 0,
\]
其中求和遍历所有与 $v$ 相邻的顶点 $w$。

施加约束 $\sum_{m=1}^{M} g_m(v) = 0$ 对所有 $v \in \mathcal{V}$ 使得问题有唯一解。
该约束恰好对应于 MILP 中的约束条件。
\end{theorem}

\begin{proof}
定义离散泛函：
\[
    \mathcal{I}_{\text{disc}}[g] := \sum_{e=(v \to w) \in \mathcal{E}}
        \sum_{m=1}^{M} \|A_{m,e} - (g_m(w) - g_m(v))\|^2.
\]

对每个 $g_m(v)$ 求偏导（$v \in \mathcal{V}, m \in \{1,\dots,M\}$）：
\[
    \frac{\partial \mathcal{I}_{\text{disc}}}{\partial g_m(v)}
    = -2 \sum_{w: (v \to w) \in \mathcal{E}} (A_{m,(v \to w)} - (g_m(w) - g_m(v)))
      + 2 \sum_{u: (u \to v) \in \mathcal{E}} (A_{m,(u \to v)} - (g_m(v) - g_m(u))).
\]

令偏导为零得到离散 Poisson 方程（图上 Laplacian）：
\[
    \sum_{w \sim v} (A_{m,(v \to w)} - (g_m(w) - g_m(v))) = 0,
\]
其中求和时边的方向性已经合适地处理。

该方程对每个 $m$ 的解存在一个加法常数自由度（若图为连通的）：
若 $g_m$ 是一个解，则 $g_m + c_m$（$c_m \in \R$ 为常数）也是解。
约束 $\sum_{m=1}^{M} g_m(v) = 0$（对所有 $v$）消除了该零模，
并且是所有此类约束中使得目标函数最优值最小的选择
（因为它是使零模投影到正交补上的约束）。
\end{proof}

\noindent
\medskip
\noindent
{\bf Theorem~\ref{thm:min_norm_discrete} (Connection Norm Minimization — Discrete Version).}
Let $\mathcal{G} = (\mathcal{V}, \mathcal{E})$ be a discretization graph of the base
manifold $\base$. Define the discrete connection potential $A: \mathcal{E} \to \R^{Md}$
as an assignment on edges. Consider the discrete variational problem:
$\min_{g: \mathcal{V} \to \R^{Md}} \sum_{e=(v \to w) \in \mathcal{E}}
\|A_e - (g(w) - g(v))\|^2$.

The stationarity condition is (for each vertex $v$ and expert $m$):
$\sum_{w \sim v} (A_{m,(v \to w)} - (g_m(w) - g_m(v))) = 0$,
where the sum runs over all vertices $w$ adjacent to $v$.

Imposing the constraint $\sum_{m=1}^{M} g_m(v) = 0$ for all $v \in \mathcal{V}$
makes the problem have a unique solution. This constraint corresponds exactly to
the constraint in the MILP formulation.

\subsection{小结：MILP 约束的必然性 / Summary: Necessity of the MILP Constraint}

\begin{corollary}[$\sum g_m = 0$ 的最优性]
约束 $\sum_{m=1}^{M} g_m = 0$ 是以下两个规范固定要求的自然融合：
\begin{enumerate}[label=(\arabic*)]
    \item 最小化联络的 $L^2$ 范数（Coulomb 规范的变分本质）；
    \item 消除整体平移零模以确保解的唯—性。
\end{enumerate}
因此，MILP 约束不是任意的，而是微分几何中最优规范固定条件的离散精确实现。
\end{corollary}

\noindent
The constraint $\sum_{m=1}^{M} g_m = 0$ is the natural fusion of two gauge-fixing requirements:
\begin{enumerate}[label=(\arabic*)]
    \item Minimizing the $L^2$ norm of the connection (the variational essence of Coulomb gauge);
    \item Eliminating the overall translation zero-mode to ensure solution uniqueness.
\end{enumerate}
Therefore, the MILP constraint is not arbitrary but the discrete exact realization
of the optimal gauge-fixing condition from differential geometry.

% ══════════════════════════════════════════════════════════════════════
\section{规范理论的代数拓扑视角 / Algebraic Topology Perspective on the Gauge Theory}
% ══════════════════════════════════════════════════════════════════════

\subsection{分类空间与万有丛 / Classifying Spaces and Universal Bundles}

在代数拓扑中，主 $G$-丛的同构类由分类空间 $BG$ 的同伦类 $[\base, BG]$ 分类。
对 $G = \R^{Md}$（可缩拓扑群），分类空间 $BG$ 是弱可缩的，故
$[\base, B\R^{Md}] = \{*\}$——所有 $\R^{Md}$-主丛都是平凡的。

然而，这并不意味所有{\bf 联络}都是平坦的。
丛的拓扑平凡性与联络的几何弯曲是两个独立的概念。
SCX 中的规范不一致性不是源于丛的拓扑非平凡性，
而是源于联络的几何弯曲——即各专家规范随参数变化的``扭曲''。

\medskip
\noindent
In algebraic topology, isomorphism classes of principal $G$-bundles are classified
by homotopy classes $[\base, BG]$ of the classifying space. For $G = \R^{Md}$
(a contractible topological group), the classifying space $BG$ is weakly contractible,
so $[\base, B\R^{Md}] = \{*\}$ — all $\R^{Md}$-principal bundles are trivial.

However, this does not mean that all {\bf connections} are flat.
The topological triviality of a bundle and the geometric curvature of a connection
are two independent concepts. The gauge inconsistency in SCX arises not from
topological non-triviality of the bundle, but from geometric curvature of the
connection — i.e., the ``twisting'' of expert gauges as parameters vary.

\subsection{联络的模空间 / Moduli Space of Connections}

\begin{definition}[联络模空间]
SCX 主丛 $\bun$ 上所有联络的空间记为 $\mathcal{A}(\bun)$。
它是仿射空间，其向量空间部分为 $\Omega^1(\base; \Ad(\bun))$
（对 Abel 群，即 $\Omega^1(\base; \R^{Md})$）。

规范群 $\mathcal{G} := C^\infty(\base, G)$ 作用于 $\mathcal{A}(\bun)$：
\[
    g \cdot \conn = \conn - \dif g \quad (\text{Abel 情形}).
\]

联络的模空间（规范轨道空间）为：
\[
    \mathcal{M} := \mathcal{A}(\bun) / \mathcal{G}.
\]
\end{definition}

\noindent
The space of all connections on the SCX principal bundle $\bun$ is denoted
$\mathcal{A}(\bun)$. It is an affine space with vector part
$\Omega^1(\base; \Ad(\bun))$ (for abelian $G$, $\Omega^1(\base; \R^{Md})$).

The gauge group $\mathcal{G} := C^\infty(\base, G)$ acts on $\mathcal{A}(\bun)$:
$g \cdot \conn = \conn - \dif g$ (abelian case).

The moduli space of connections (gauge orbit space) is:
$\mathcal{M} := \mathcal{A}(\bun) / \mathcal{G}$.

\begin{proposition}[Cercis 作为模空间上的函数]
Cercis 分数 $\cercis: \mathcal{A}(\bun) \to \R_{\geq 0}$ 在规范轨道上是常数，
故它下降为模空间上的良定义函数：
\[
    \cercis: \mathcal{M} \to \R_{\geq 0}.
\]
该函数的零点集 $\cercis^{-1}(0)$ 恰好是平坦联络的模空间
$\mathcal{M}_{\text{flat}} \subset \mathcal{M}$。
\end{proposition}

\noindent
The Cercis Score $\cercis: \mathcal{A}(\bun) \to \R_{\geq 0}$ is constant on gauge orbits,
hence descends to a well-defined function on the moduli space:
$\cercis: \mathcal{M} \to \R_{\geq 0}$.
Its zero set $\cercis^{-1}(0)$ is precisely the moduli space of flat connections
$\mathcal{M}_{\text{flat}} \subset \mathcal{M}$.

\subsection{瞬子与最优化 / Instantons and Optimization}

连续 Yang-Mills 理论中，Yang-Mills 泛函 $S_{\text{YM}}[A] = \int \|F\|^2 \vol$
的极小值（即 Yang-Mills 瞬子）满足自对偶方程
$F = \pm \star F$（在 4 维情形）。

在 SCX 中，Cercis 泛函的极小值给出了多专家系统的最优规范调整方案，
其在离散情形的解就是 MILP 的解。

\begin{proposition}[MILP 作为最优化问题]
MILP 的目标函数是离散 Cercis 泛函（加上规范固定约束）：
\[
    \min_{A \in \mathcal{A}_{\text{disc}}} \cercis_{\text{disc}}(A)
    \quad \text{s.t.} \quad \sum_m g_m = 0.
\]
该问题的最优点 $A^*$ 是离散联络模空间中满足``库仑条件''的联络。
\end{proposition}

\noindent
In continuous Yang-Mills theory, minima of the Yang-Mills functional
$S_{\text{YM}}[A] = \int \|F\|^2 \vol$ (i.e., Yang-Mills instantons) satisfy the
self-duality equation $F = \pm \star F$ (in 4 dimensions).

In SCX, minima of the Cercis functional give the optimal gauge adjustment scheme
for the multi-expert system, whose discrete solution is precisely the MILP solution.

The MILP objective function is the discrete Cercis functional (plus gauge-fixing constraint):
$\min_{A \in \mathcal{A}_{\text{disc}}} \cercis_{\text{disc}}(A)$
s.t.\ $\sum_m g_m = 0$.
The optimal point $A^*$ of this problem is the connection in the discrete moduli space
satisfying the ``Coulomb condition.''

% ══════════════════════════════════════════════════════════════════════
\section{算法实现与数值考量 / Algorithmic Implementation and Numerical Considerations}
% ══════════════════════════════════════════════════════════════════════

\subsection{离散联络的构造 / Construction of the Discrete Connection}

给定 $M$ 个专家在 $N$ 个参数构型 $\{x^k\}_{k=1}^N \subset \mathcal{X}$ 上的预测，
SCX 联络的构造如下：

\begin{enumerate}[label=步骤 \arabic*:, leftmargin=*]
    \item {\bf 成对专家位移:} 对每个参数构型 $k$ 和每对专家 $(i, j)$，
        计算原始坐标差 $\delta_{ij}^k := \tilde{x}_i^k - \tilde{x}_j^k$。
        此即图 $\mathcal{G}$ 上的边赋值 $A_e$。
    \item {\bf 构建联络图:} 构造有向图 $\mathcal{G} = (\mathcal{V}, \mathcal{E})$，
        其中 $\mathcal{V} = \{1, \dots, N\} \times \{1, \dots, M\}$，
        $\mathcal{E}$ 包含两类边：
        \begin{itemize}
            \item 参数边: $(k, m) \to (k+1, m)$，表示同一专家沿参数轴的移动；
            \item 专家边: $(k, i) \to (k, j)$，表示同一参数点处不同专家的比较。
        \end{itemize}
    \item {\bf 曲率计算:} 对图中每个基本回路（四边形），计算离散和乐：
        \[
            \curv_{(k,i,j)} = \delta_{ij}^k + \delta_{ji}^{k+1} + \delta_{ij}^{k+1} + \delta_{ji}^k \quad (\text{循环和}).
        \]
\end{enumerate}

\medskip
\noindent
Given predictions of $M$ experts on $N$ parameter configurations
$\{x^k\}_{k=1}^N \subset \mathcal{X}$, the SCX connection is constructed as follows:

\begin{enumerate}[label=Step \arabic*:, leftmargin=*]
    \item {\bf Pairwise expert displacements:} For each parameter configuration $k$
        and each expert pair $(i, j)$, compute the raw coordinate difference
        $\delta_{ij}^k := \tilde{x}_i^k - \tilde{x}_j^k$.
        This serves as the edge assignment $A_e$ on the graph $\mathcal{G}$.
    \item {\bf Build connection graph:} Construct a directed graph
        $\mathcal{G} = (\mathcal{V}, \mathcal{E})$, where
        $\mathcal{V} = \{1, \dots, N\} \times \{1, \dots, M\}$,
        and $\mathcal{E}$ contains two types of edges:
        \begin{itemize}
            \item Parameter edges: $(k, m) \to (k+1, m)$, representing the movement
                of the same expert along the parameter axis;
            \item Expert edges: $(k, i) \to (k, j)$, representing the comparison
                of different experts at the same parameter point.
        \end{itemize}
    \item {\bf Curvature computation:} For each elementary loop (square) in the graph,
        compute the discrete holonomy:
        $\curv_{(k,i,j)} = \delta_{ij}^k + \delta_{ji}^{k+1} + \delta_{ij}^{k+1} + \delta_{ji}^k$
        (cyclic sum).
\end{enumerate}

\subsection{Cercis 计算与 MILP 求解 / Cercis Computation and MILP Solving}

\begin{algorithm}[H]
\caption{SCX Cercis 与 MILP 计算}
\begin{enumerate}[label=(\arabic*)]
    \item 输入：$M$ 个专家在 $N$ 个构型上的预测 $\{\tilde{x}_m^k\}$。
    \item 计算成对位移矩阵 $\Delta^k = (\delta_{ij}^k)_{i,j=1}^M$ 对所有 $k$。
    \item 计算离散曲率（对所有四边形回路）：
        \[
            \curv_{(k,i,j)} = \delta_{ij}^k + \delta_{ji}^{k+1} + \delta_{ij}^{k+1} + \delta_{ji}^k.
        \]
    \item 计算 Cercis 分数：
        \[
            \cercis = \frac{1}{N M^2} \sum_{k=1}^{N-1} \sum_{i,j=1}^{M} \|\curv_{(k,i,j)}\|^2.
        \]
    \item 若 $\cercis \approx 0$：输出``系统规范一致''，全局对齐直接可用。
    \item 若 $\cercis > \tau$（阈值 $\tau$ 由数值容差确定）：
        \begin{enumerate}
            \item 构建 MILP 问题，目标函数为离散 $L^2$ 范数；
            \item 施加约束 $\sum_m g_m^k = 0$；
            \item 求解 MILP 得到最优规范调整 $\{g_m^k\}$；
            \item 输出规范对齐后的预测 $\hat{x}_m^k = \tilde{x}_m^k - g_m^k$。
        \end{enumerate}
\end{enumerate}
\end{algorithm}

\subsection{收敛性分析 / Convergence Analysis}

\begin{theorem}[收敛性]
设 $\base$ 为紧致 Riemann 流形，$\{\base_h\}_{h > 0}$ 为其正则三角剖分序列
（$h$ 为网格尺寸）。设 $A_h$ 为 $\base_h$ 上的离散联络。
则当 $h \to 0$ 时：
\begin{enumerate}[label=(\roman*)]
    \item $\cercis_{\text{disc}}(A_h) \to \cercis(A)$（逐点收敛）；
    \item MILP 的解 $g_h$ 弱收敛于连续 Coulomb 规范固定问题的解 $g$。
\end{enumerate}
\end{theorem}

\noindent
Let $\base$ be a compact Riemannian manifold and $\{\base_h\}_{h > 0}$ a sequence
of regular triangulations ($h$ is the mesh size). Let $A_h$ be the discrete connection
on $\base_h$. Then as $h \to 0$:
(i) $\cercis_{\text{disc}}(A_h) \to \cercis(A)$ (pointwise convergence);
(ii) The MILP solution $g_h$ converges weakly to the solution $g$ of the continuous
Coulomb gauge-fixing problem.

% ══════════════════════════════════════════════════════════════════════
\section{泛化与推广 / Generalizations and Extensions}
% ══════════════════════════════════════════════════════════════════════

\subsection{非 Abel 专家群 / Non-Abelian Expert Groups}

本文讨论的 SCX 基于 Abel 规范群 $G_m \cong (\R^d, +)$，对应简单的平移对称性。
若专家的规范群扩展为包含旋转、缩放等变换，则 $G_m$ 变为非 Abel Lie 群
（如 $\mathrm{SE}(d)$、$\mathrm{GL}(d, \R)$）。
在此情况下：

\begin{itemize}
    \item 联络 $1$-形式 $\conn$ 取非 Abel 值，满足 $R_g^* \conn = \Ad_{g^{-1}} \circ \conn$；
    \item 曲率 $\curv = \dif \conn + \frac{1}{2}[\conn, \conn]$ 包含非线性项；
    \item 规范变换 $A' = g^{-1} A g + g^{-1} \dif g$ 在非 Abel 情形更为复杂；
    \item Cercis 分数 $\cercis = \int \Tr(F \wedge \star F)$ 保持规范不变；
    \item MILP 需修改为混合整数非线性规划 (MINLP)。
\end{itemize}

这些推广对应于更一般的``专家对称性''，本文的框架提供了自然的推广路径。

\medskip
\noindent
This paper discusses SCX based on the abelian gauge group $G_m \cong (\R^d, +)$,
corresponding to simple translational symmetry.
If the gauge group of experts is extended to include rotations, scaling, etc.,
then $G_m$ becomes a non-abelian Lie group (e.g., $\mathrm{SE}(d)$, $\mathrm{GL}(d, \R)$).
In this case:
\begin{itemize}
    \item The connection $1$-form $\conn$ takes non-abelian values, satisfying
        $R_g^* \conn = \Ad_{g^{-1}} \circ \conn$;
    \item The curvature $\curv = \dif \conn + \frac{1}{2}[\conn, \conn]$ includes nonlinear terms;
    \item Gauge transformations $A' = g^{-1} A g + g^{-1} \dif g$ are more complex;
    \item The Cercis Score $\cercis = \int \Tr(F \wedge \star F)$ remains gauge-invariant;
    \item MILP must be modified to Mixed-Integer Nonlinear Programming (MINLP).
\end{itemize}
These generalizations correspond to more general ``expert symmetries,''
and the framework of this paper provides a natural path for extension.

\subsection{无穷维极限与连续专家场 / Infinite-Dimensional Limit and Continuous Expert Fields}

当专家数量 $M \to \infty$ 时，离散专家指标 $m$ 趋向于连续参数 $s \in [0,1]$。
此时规范群 $G = \prod_m G_m$ 趋向于连续映射群
$G_\infty = C^\infty([0,1], \R^d)$ 或其适当的 Sobolev 完备化。

在此极限下，主丛 $\bun$ 变为无穷维丛，
联络 $\conn(s, x)$ 成为两个变量的函数，
曲率 $\curv = \dif \conn$ 包含两个方向：
沿参数空间的方向（$\dif_x$）和沿专家指标的方向（$\dif_s$）。

Cercis 分数变为双积分：
\[
    \cercis_\infty = \int_0^1 \int_{\base} \|F(s, x)\|^2 \, \vol_{\base} \, \dif s.
\]

这开启了一个有趣的研究方向：SCX 理论与二维规范场论的联系。

\medskip
\noindent
As the number of experts $M \to \infty$, the discrete expert index $m$ tends toward
a continuous parameter $s \in [0,1]$. The gauge group $G = \prod_m G_m$ then tends
to the continuous mapping group $G_\infty = C^\infty([0,1], \R^d)$ or its appropriate
Sobolev completion.

In this limit, the principal bundle $\bun$ becomes an infinite-dimensional bundle,
the connection $\conn(s, x)$ becomes a function of two variables,
and the curvature $\curv = \dif \conn$ includes two directions:
the parameter-space direction ($\dif_x$) and the expert-index direction ($\dif_s$).

The Cercis Score becomes a double integral:
$\cercis_\infty = \int_0^1 \int_{\base} \|F(s, x)\|^2 \, \vol_{\base} \, \dif s$.

This opens an interesting research direction: the connection between SCX theory
and two-dimensional gauge field theory.

\subsection{与拓扑数据分析的联系 / Connection to Topological Data Analysis}

SCX 曲率 $\curv$（作为非零上同调类）与拓扑数据分析中的持续同调 (persistent homology)
有自然的联系：
\begin{itemize}
    \item 非零离散曲率回路对应于图 $\mathcal{G}$ 中的 $1$-维``洞''；
    \item Cercis 分数可以解释为持续 $1$-维同调类在某种过滤中的持续性加权和；
    \item 曲率的局部化（通过 MILP）等价于寻找极小化``$1$-维空洞''的最优上边缘。
\end{itemize}

\medskip
\noindent
The SCX curvature $\curv$ (as a non-zero cohomology class) has natural connections
to persistent homology in topological data analysis:
\begin{itemize}
    \item Non-zero discrete curvature loops correspond to $1$-dimensional ``holes''
        in the graph $\mathcal{G}$;
    \item The Cercis Score can be interpreted as the persistence-weighted sum of
        persistent $1$-dimensional homology classes under some filtration;
    \item Localizing curvature (via MILP) is equivalent to finding the optimal
        coboundary that minimizes ``$1$-dimensional voids.''
\end{itemize}

% ══════════════════════════════════════════════════════════════════════
\section{结论与展望 / Conclusion and Outlook}
% ══════════════════════════════════════════════════════════════════════

\subsection{主要结果总结 / Summary of Main Results}

本文从微分几何的主纤维丛理论出发，严格建立了 SCX 多专家系统的规范理论形式化。
主要结论如下：

\begin{enumerate}[label=(\arabic*)]
    \item {\bf SCX 规范群：} 每个专家的平移群 $G_m \cong (\R^d, +)$ 构成 Abel 规范群，
        总规范群为 $G = \prod_{m=1}^M G_m \cong \R^{Md}$。

    \item {\bf 主纤维丛结构：} 全空间 $\bun = \R^{Md} \times \mathcal{X}$，
        底流形 $\base = \bun / G \cong \mathcal{X}$，
        投影 $\pi: \bun \to \base$ 构成主 $G$-丛。

    \item {\bf 联络与曲率：} 联络 $1$-形式 $\conn \in \Omega^1(\bun, \g)$ 描述
        专家规范随参数的变化；曲率 $\curv = \dif \conn$ 度量平行移动的路径依赖性。

    \item {\bf 全局规范固定定理：} $\curv = 0$ 当且仅当存在全局规范固定
        （所有专家共享同一规范），这是由平坦联络的 Frobenius 可积性保证的。

    \item {\bf 势能面不齐定理：} $\curv \neq 0$ 时，不存在同时对齐所有成对专家比较的
        单一规范选择——这是势能面不齐的几何根源。

    \item {\bf Cercis 分数：} $\cercis = \int \|F\|^2 \vol$ 是规范不变的几何可观测量，
        类似于 Yang-Mills 理论中的 $\Tr(F_{\mu\nu}F^{\mu\nu})$。
        $\cercis = 0$ 判定系统的规范一致性，$\cercis > 0$ 量化不一致程度。

    \item {\bf MILP 作为 Coulomb 规范：} MILP 规范固定是离散的 Coulomb 规范选择，
        其约束 $\sum_m g_m = 0$ 最小化联络的 $L^2$ 范数且消除整体平移零模。

    \item {\bf 变分定理：} $\sum_m g_m = 0$ 约束是离散 Coulomb 规范的 Euler-Lagrange
        方程的自然结果，确保了规范固定方案的唯一性和最优性。
\end{enumerate}

\medskip
\noindent
This paper rigorously establishes the gauge-theoretic formalization of SCX multi-expert
systems using the theory of principal fiber bundles in differential geometry.
The main results are:
\begin{enumerate}[label=(\arabic*)]
    \item {\bf SCX Gauge Group:} Each expert's translation group $G_m \cong (\R^d, +)$
        forms an abelian gauge group; the total gauge group is
        $G = \prod_{m=1}^M G_m \cong \R^{Md}$.
    \item {\bf Principal Bundle Structure:} Total space $\bun = \R^{Md} \times \mathcal{X}$,
        base manifold $\base = \bun / G \cong \mathcal{X}$,
        projection $\pi: \bun \to \base$ constitutes a principal $G$-bundle.
    \item {\bf Connection and Curvature:} The connection $1$-form
        $\conn \in \Omega^1(\bun, \g)$ describes how expert gauges vary with parameters;
        curvature $\curv = \dif \conn$ measures path-dependence of parallel transport.
    \item {\bf Global Gauge-Fixing Theorem:} $\curv = 0$ iff global gauge-fixing exists
        (all experts share the same gauge), guaranteed by Frobenius integrability
        of flat connections.
    \item {\bf Potential Energy Surface Misalignment Theorem:} When $\curv \neq 0$,
        no single gauge choice can simultaneously align all pairwise expert comparisons —
        this is the geometric origin of PES misalignment.
    \item {\bf Cercis Score:} $\cercis = \int \|F\|^2 \vol$ is a gauge-invariant
        geometric observable, analogous to $\Tr(F_{\mu\nu}F^{\mu\nu})$ in Yang-Mills theory.
        $\cercis = 0$ determines gauge consistency; $\cercis > 0$ quantifies inconsistency.
    \item {\bf MILP as Coulomb Gauge:} MILP gauge-fixing is a discrete Coulomb gauge
        selection; its constraint $\sum_m g_m = 0$ minimizes the $L^2$ norm of the
        connection and eliminates the overall translation zero-mode.
    \item {\bf Variational Theorem:} The $\sum_m g_m = 0$ constraint is a natural consequence
        of the Euler-Lagrange equation for the discrete Coulomb gauge, ensuring uniqueness
        and optimality of the gauge-fixing scheme.
\end{enumerate}

\subsection{展望 / Outlook}

本工作开启了若干值得进一步探索的方向：
\begin{enumerate}[label=(\arabic*)]
    \item {\bf 非 Abel 推广：} 将规范群扩展至 $\mathrm{SE}(d)$、$\mathrm{GL}(d)$ 等，
        以包含旋转和缩放对称性。
    \item {\bf 瞬子解与最优化：} 研究 Cercis 泛函的瞬子解（自对偶联络），
        探索其与最优传输理论（optimal transport）和 Wasserstein 几何的联系。
    \item {\bf Hodge 理论与调和形式：} 利用 Hodge 分解理论深入分析
        联络空间的几何结构，构造更一般的规范不变量。
    \item {\bf 与机器学习理论的联系：} 将曲率解释为模型泛化误差的几何度量，
        建立 SCX 规范理论与学习理论的桥梁。
    \item {\bf 数值算法优化：} 基于本文的几何洞察，设计更高效的离散规范固定算法，
        可能利用谱方法、多重网格方法或黎曼优化。
    \item {\bf 无穷维极限：} 严格研究 $M \to \infty$ 极限下的连续专家场理论
        及其与二维规范场论的关系。
\end{enumerate}

\medskip
\noindent
This work opens several directions worthy of further exploration:
\begin{enumerate}[label=(\arabic*)]
    \item {\bf Non-abelian generalization:} Extending the gauge group to $\mathrm{SE}(d)$,
        $\mathrm{GL}(d)$, etc., to include rotational and scaling symmetries.
    \item {\bf Instanton solutions and optimization:} Studying instanton solutions
        (self-dual connections) of the Cercis functional and exploring connections
        with optimal transport theory and Wasserstein geometry.
    \item {\bf Hodge theory and harmonic forms:} Using Hodge decomposition to deeply
        analyze the geometry of connection space and construct more general gauge invariants.
    \item {\bf Connections to machine learning theory:} Interpreting curvature as a geometric
        measure of model generalization error, building a bridge between SCX gauge theory
        and learning theory.
    \item {\bf Numerical algorithm optimization:} Based on the geometric insights of this paper,
        designing more efficient discrete gauge-fixing algorithms, potentially using
        spectral methods, multigrid methods, or Riemannian optimization.
    \item {\bf Infinite-dimensional limit:} Rigorously studying the continuous expert field
        theory in the $M \to \infty$ limit and its relationship to two-dimensional gauge theory.
\end{enumerate}

% ══════════════════════════════════════════════════════════════════════
% 参考文献 / References
% ══════════════════════════════════════════════════════════════════════

\begin{thebibliography}{99}

\bibitem{kobayashi1963}
S.~Kobayashi and K.~Nomizu.
\emph{Foundations of Differential Geometry, Vol.~I}.
Wiley-Interscience, 1963.

\bibitem{ehresmann1950}
C.~Ehresmann.
``Les connexions infinitésimales dans un espace fibré différentiable.''
\emph{Colloque de Topologie, Bruxelles}, 29--55, 1950.

\bibitem{weyl1918}
H.~Weyl.
``Gravitation und Elektrizität.''
\emph{Sitzungsberichte der Königlich Preußischen Akademie der Wissenschaften},
465--480, 1918.

\bibitem{yangmills1954}
C.~N.~Yang and R.~L.~Mills.
``Conservation of Isotopic Spin and Isotopic Gauge Invariance.''
\emph{Physical Review}, 96(1):191--195, 1954.

\bibitem{atiyah1978}
M.~F.~Atiyah, N.~J.~Hitchin, and I.~M.~Singer.
``Self-duality in four-dimensional Riemannian geometry.''
\emph{Proceedings of the Royal Society A}, 362(1711):425--461, 1978.

\bibitem{donaldson1983}
S.~K.~Donaldson.
``An application of gauge theory to four-dimensional topology.''
\emph{Journal of Differential Geometry}, 18(2):279--315, 1983.

\bibitem{witten1988}
E.~Witten.
``Topological quantum field theory.''
\emph{Communications in Mathematical Physics}, 117(3):353--386, 1988.

\bibitem{nahm1980}
W.~Nahm.
``A simple formalism for the BPS monopole.''
\emph{Physics Letters B}, 90(4):413--414, 1980.

\bibitem{bleecker1981}
D.~Bleecker.
\emph{Gauge Theory and Variational Principles}.
Addison-Wesley, 1981.

\bibitem{freed1995}
D.~S.~Freed and K.~K.~Uhlenbeck.
\emph{Instantons and Four-Manifolds}, 2nd ed.
Springer, 1995.

\bibitem{husemoller1966}
D.~Husemoller.
\emph{Fibre Bundles}.
McGraw-Hill, 1966 (3rd ed., Springer, 1994).

\bibitem{spanier1981}
E.~H.~Spanier.
\emph{Algebraic Topology}.
Springer, 1981.

\bibitem{nakahara2003}
M.~Nakahara.
\emph{Geometry, Topology and Physics}, 2nd ed.
Taylor \& Francis, 2003.

\bibitem{guillemin2013}
V.~Guillemin and S.~Sternberg.
\emph{Symplectic Techniques in Physics}, 2nd ed.
Cambridge University Press, 2013.

\bibitem{hatcher2002}
A.~Hatcher.
\emph{Algebraic Topology}.
Cambridge University Press, 2002.

\bibitem{bott1982}
R.~Bott and L.~Tu.
\emph{Differential Forms in Algebraic Topology}.
Springer, 1982.

\end{thebibliography}

\end{document}
